\documentclass[11pt,reqno]{amsart}
\usepackage[T1]{fontenc}\usepackage[utf8]{inputenc}\usepackage{lmodern}
\usepackage{amsmath,amssymb,amsfonts,amsthm}\usepackage[a4paper,margin=28mm]{geometry}
\usepackage{booktabs,array,longtable,graphicx}\usepackage[colorlinks=true,allcolors=blue]{hyperref}
\setlength{\emergencystretch}{4em}\raggedbottom\numberwithin{equation}{section}
\theoremstyle{plain}\newtheorem{theorem}{Theorem}[section]\newtheorem{proposition}[theorem]{Proposition}\newtheorem{lemma}[theorem]{Lemma}\newtheorem{corollary}[theorem]{Corollary}
\theoremstyle{definition}\newtheorem{hypothesis}[theorem]{Hypothesis}\newtheorem{definition}[theorem]{Definition}
\theoremstyle{remark}\newtheorem{remark}[theorem]{Remark}\newtheorem{experiment}[theorem]{Experiment}\newtheorem{question}[theorem]{Question}\newtheorem{heuristic}[theorem]{Heuristic}\newtheorem{conjecture}[theorem]{Conjecture}
\newcommand{\NO}{N_0}\newcommand{\lL}{\mathcal L}\newcommand{\IT}{\mathcal I_T}\newcommand{\R}{\mathbb R}\newcommand{\N}{\mathbb N}\newcommand{\eps}{\varepsilon}\newcommand{\sgn}{\operatorname{sgn}}
\newcommand{\vpi}{\varpi}\newcommand{\mathscr}[1]{\mathcal{#1}}\newcommand{\Z}{\mathbb Z}\newcommand{\C}{\mathbb C}\DeclareMathOperator{\Li}{Li}
\title[Ordinal defects and prime cutoffs]{Ordinal defects, variable prime cutoffs, and response identities for Riemann zeta zero ordinates}
\author{V\'ictor Casa\~na-Giner}\address{Valencia, Spain}\email{Email to be inserted}\subjclass[2020]{11M26, 11M06, 11Y35}\keywords{Riemann zeta-function, zero ordinates, Selberg approximation, Euler phases, pair correlation}\date{17 October 2026}
\begin{document}
\begin{abstract}
Let $\gamma_k\in\R_{>0}$ be the positive ordinates of the nontrivial zeros of $\zeta(s)$, repeated according to multiplicity.  We study defects obtained by comparing the ordinal $k-1/2$ with the smooth Riemann--von Mangoldt count and finite Euler phases.  The paper establishes a quadrature identity at the jumps of a counting function; moment formulae for the midpoint-clock displacement; a uniform mean-square estimate for a cutoff $2\le X\le T^{1/2}$; localization under a strict derivative condition; fixed-frequency projection and finite-window least-squares identities; a decomposition of the one-step clock; a conditional moving-prime response; a distributional Fourier representation of its scalar target; and an evaluation in a finite determinantal model.  Every theorem states the domains of its variables, and unconditional, conditional, and model statements are separated.  Finite computations are reported only for observables defined in the text.
\end{abstract}\maketitle

\section{Conventions and logical status}\label{sec:conventions}
The symbol $k$ always denotes a positive integer.  The variables $t,T,X,x_0$ are real unless explicitly declared otherwise; asymptotic statements use $T\to\infty$ through real values.  A prime variable $p$ ranges over rational primes, and an integer frequency $n$ lies in $\mathbb Z_{\ge2}$.  Complex zeros are denoted $\rho=\beta+i\gamma\in\mathbb C$, while $\gamma\in\mathbb R$ denotes an ordinate.  Sums over zeros count multiplicity.  All integrals are Lebesgue integrals; on the piecewise continuous functions occurring below they equal the corresponding improper Riemann integrals.  Big-$O$ and $\ll$ constants are absolute unless their dependence is displayed by a subscript.

Sections~\ref{sec:identities}--\ref{sec:one_step_clock} contain unconditional statements except where a named hypothesis is invoked.  Section~\ref{sec:moving_response} is conditional.  Section~\ref{sec:Fourier_reduction} is an unconditional distributional reduction.  Section~\ref{sec:sine_response} concerns a specified finite determinantal model and is not asserted for zeta zeros.
\section{Notation and the clock}
\label{sec:notation}

\subsection{Counting functions}
\label{subsec:counting}
Let $\vartheta(t)=\Im\log\Gamma(\frac14+\frac{it}2)-\frac t2\log\pi$ be the Riemann--Siegel theta function. For $t>0$ not an ordinate \cite[Thms.~9.3--9.4]{Titchmarsh1986},
\begin{equation}
N(t)=\frac1\pi\vartheta(t)+1+S(t),\qquad S(t)=\frac1\pi\arg\zeta\bigl(\tfrac12+it\bigr)=O(\log t),
\label{eq:RvM}
\end{equation}
with $\arg\zeta$ defined by continuous variation from $2$ along the segment to $2+it$ and then horizontally to $\frac12+it$. We set
\begin{equation}
\Theta_0(t)=\frac t2\log\frac t{2\pi e}-\frac\pi8,\qquad \NO(t)=\frac1\pi\Theta_0(t)+1,\qquad \omega(t)=\vartheta(t)-\Theta_0(t),
\label{eq:N0_def}
\end{equation}
so that $\omega(t)=\frac1{48t}+O(t^{-3})$ by Stirling's formula and
\begin{equation}
N(t)=\NO(t)+S(t)+\frac{\omega(t)}\pi .
\label{eq:N_decomposition}
\end{equation}
The derivative of the smooth count is
\begin{equation}
\lL(t)=\NO'(t)=\frac1{2\pi}\log\frac t{2\pi}.
\label{eq:L_def}
\end{equation}
For $T\ge 1$ let $\IT=\{k\ge1:T<\gamma_k\le2T\}$. Since $N(2T)-N(T)=\NO(2T)-\NO(T)+O(\log T)$ by \eqref{eq:RvM},
\begin{equation}
\#\IT=T\lL(T)+O(T).
\label{eq:IT_size}
\end{equation}

\subsection{The clock and the displacements}
\label{subsec:clock}

\begin{equation}\label{eq:clock}
\tau_k=\frac{2\pi(k-11/8)}{W_0((k-11/8)/e)},\qquad k\in\mathbb Z_{\ge2},
\end{equation}
\begin{lemma}\label{lem:clock}
The sequence \eqref{eq:clock} is the unique solution of $\NO(\tau_k)=k-\frac12$ with $\tau_k>2\pi$, $k\ge2$. As a function of a real variable $k$, $\tau_k$ is smooth and increasing with
\begin{equation}
\frac{d\tau_k}{dk}=\frac1{\lL(\tau_k)}=\frac{2\pi}{\log(\tau_k/2\pi)}.
\label{eq:clock_derivative}
\end{equation}
\end{lemma}

\begin{proof}
With $x=t/2\pi$ the equation $\NO(t)=k-\frac12$ reads $x\log(x/e)=k-\frac{11}8$. Putting $x=ey$ gives $y\log y=(k-\frac{11}8)/e$, so $\log y=W((k-\frac{11}8)/e)$ with the principal branch \cite{Corless1996} and $x=(k-\frac{11}8)/W((k-\frac{11}8)/e)$. The map $x\mapsto x\log(x/e)$ is increasing for $x>1$, which gives uniqueness, and \eqref{eq:clock_derivative} follows by differentiating $\NO(\tau_k)=k-\frac12$.
\end{proof}

\begin{definition}\label{def:displacements}
The ordinal displacement and the displacement of the clock are
\begin{equation}
s_k=k-\frac12-\NO(\gamma_k),\qquad \delta_k=\tau_k-\gamma_k .
\label{eq:displacements}
\end{equation}
\end{definition}

By Lemma~\ref{lem:clock}, $s_k=\NO(\tau_k)-\NO(\gamma_k)$, and the mean value theorem gives
\begin{equation}
s_k=\lL(\xi_k)\,\delta_k\qquad\text{for some }\xi_k\text{ between }\gamma_k\text{ and }\tau_k .
\label{eq:s_delta}
\end{equation}

\begin{remark}\label{rem:gram_points}
Since $\Theta_0(\tau_k)=\pi(k-\frac32)$ and $\vartheta=\Theta_0+\omega$, the clock point $\tau_k$ is the point at which $\vartheta$ takes the value $\pi(k-\frac32)$, up to a shift $O(\tau_k^{-1}(\log\tau_k)^{-1})$. Thus $\tau_k$ lies, in phase, half way between the Gram points $g_{k-2}$ and $g_{k-1}$, defined by $\vartheta(g_n)=n\pi$. The integer displacement of $\gamma_k$ with respect to the Gram intervals is the object of Gram's law; see Remark~\ref{rem:korolev}.
\end{remark}

\begin{lemma}\label{lem:s_bounds}
For every $k$ with $\gamma_k>2$,
\begin{equation}
N(\gamma_k^-)-\NO(\gamma_k)+\frac12\le s_k\le N(\gamma_k^+)-\NO(\gamma_k)-\frac12,
\label{eq:s_bounds}
\end{equation}
where $N(\gamma^\pm)$ are the one-sided limits. In particular $|s_k|\le|S(\gamma_k^-)|+|S(\gamma_k^+)|+1$.
\end{lemma}

\begin{proof}
The indices $j$ with $\gamma_j=\gamma_k$ are $N(\gamma_k^-)+1,\dots,N(\gamma_k^+)$, hence $N(\gamma_k^-)+1\le k\le N(\gamma_k^+)$; subtract $\frac12+\NO(\gamma_k)$ and use \eqref{eq:N_decomposition} with $|\omega|/\pi<\frac12$.
\end{proof}

\subsection{Euler phases and the truncated equation}
\label{subsec:euler_phase}
For $n\ge2$ put
\begin{equation}
a_n=-\frac{\Lambda(n)}{\pi\sqrt n\,\log n},
\label{eq:an_def}
\end{equation}
where $\Lambda$ is the von Mangoldt function. For a set $A$ of prime powers with $\sum_{n\in A}|a_n|<\infty$ let
\begin{equation}
S_A(t)=\sum_{n\in A}a_n\sin(t\log n),\qquad v_A=\frac12\sum_{n\in A}a_n^2 .
\label{eq:SA_def}
\end{equation}
For $X\ge2$ let $A_X=\{p^m:p\le X,\ m\ge1\}$ and $S_X=S_{A_X}$. Since $\arg(1-z)=-\Im\sum_{m\ge1}z^m/m$ for $|z|<1$,
\begin{equation}
S_X(t)=-\frac1\pi\sum_{p\le X}\arg\bigl(1-p^{-1/2-it}\bigr)=-\frac1\pi\sum_{p\le X}\sum_{m\ge1}\frac{\sin(mt\log p)}{m\,p^{m/2}},
\label{eq:SX_def}
\end{equation}
For each fixed real $X\ge2$, the outer prime sum is finite and the inner series is absolutely and uniformly convergent for $t\in\mathbb R$, because $\sum_{m\ge1}(m p^{m/2})^{-1}<\infty$.  The differentiated series is also uniformly convergent, since $\sum_{m\ge1}(\log p)p^{-m/2}<\infty$.  Consequently $S_X:\mathbb R\to\mathbb R$ is $C^1$, and integration and differentiation term by term are justified without an interchange involving an infinite prime sum.
with the principal argument of each individual factor. The sum is a continuous lift of the argument of the product; it need not be its principal argument. For a finite set of primes $P$ we write $S_P$ for the same sum over $p\in P$, and
\begin{equation}
c_P=\frac1{2\pi^2}\sum_{p\in P}\mathrm{Li}_2\Bigl(\frac1p\Bigr)=\lim_{M\to\infty}v_{\{p^m:\,p\in P,\ m\le M\}},
\label{eq:cP_def}
\end{equation}
where $\mathrm{Li}_2(x)=\sum_{m\ge1}x^m/m^2$. We also use the bounds
\begin{equation}
|S_X(t)|\le\sigma_X:=\frac1\pi\sum_{p\le X}\sum_{m\ge1}\frac1{m\,p^{m/2}},\qquad |S_X'(t)|\le M_X:=\frac1\pi\sum_{p\le X}\frac{\log p}{\sqrt p-1},
\label{eq:sigma_M_def}
\end{equation}
and note that $\sigma_X\ll\sqrt X/\log X$ and $M_X\ll\sqrt X$ by the prime number theorem.

\begin{definition}\label{def:LeClair_model}
For $X\ge2$ the truncated counting model and the defect of the $k$-th zero are
\begin{equation}
F_X(t)=\NO(t)+S_X(t),\qquad u_k(X)=k-\frac12-F_X(\gamma_k)=s_k-S_X(\gamma_k).
\label{eq:defect_def}
\end{equation}
The one-step corrected clock is
\begin{equation}
\tau_k^X=\tau_k-\frac{S_X(\tau_k)}{\lL(\tau_k)}=\tau_k-\frac{d\tau_k}{dk}\,S_X(\tau_k),
\label{eq:one_step_clock}
\end{equation}
and, for a saturation level $B>0$, $\tau_k^{X,B}=\tau_k-\mathrm{sat}_B(S_X(\tau_k)/\lL(\tau_k))$ with $\mathrm{sat}_B(y)=\sgn(y)\min(|y|,B)$. We write $\delta_k^X=\tau_k^X-\gamma_k$ and $\delta_k^{X,B}=\tau_k^{X,B}-\gamma_k$.
\end{definition}

\begin{lemma}\label{lem:LeClair_equation}
The exact counting relation is $N(t)=\NO(t)+S(t)+\omega(t)/\pi$, with $\omega(t)=O(1/t)$, as specified in \eqref{eq:N_decomposition}. At a simple ordinate its midpoint value is
\begin{equation}
 \NO(\gamma_k)+S^*(\gamma_k)+\frac{\omega(\gamma_k)}\pi
                  =k-\frac12.
 \label{eq:exact_midpoint_ordinate}
\end{equation}
For a repeated ordinate the right side is $\bar k_\gamma-\frac12$. Replacing $S$ by $S_X$ and omitting the smooth remainder defines the model equation $F_X(t)=k-\frac12$. The point $\tau_k^X$ is the first step of the iteration $t\mapsto t-(F_X(t)-k+\frac12)/\lL(t)$ started at $t=\tau_k$.
\end{lemma}
\begin{proof}
The midpoint of the counting jump is $N(\gamma^-)+\nu(\gamma)/2=\bar k_\gamma-1/2$, so the exact count decomposition gives \eqref{eq:exact_midpoint_ordinate} and its repeated-ordinate form. The continuous argument of the finite Euler product, obtained by summing the principal arguments of its individual local factors, is $-\sum_{p\le X}\arg(1-p^{-1/2-it})=\pi S_X(t)$. Finally $F_X(\tau_k)-k+\frac12=S_X(\tau_k)$, which gives the first iterate. The arithmetic defect studied below uses the stated model; its exact identities never discard the smooth remainder in the genuine count.
\end{proof}

\section{Known results used}
\label{sec:known_results}

This section collects standard material; nothing in it is new.

\begin{theorem}[Riemann--von Mangoldt; Littlewood; {\cite[Thms.~9.4, 9.9 and \S9.4]{Titchmarsh1986}}, \cite{Littlewood1924}]\label{thm:known_basic}
$S(t)=O(\log t)$; the gaps $\gamma_{k+1}-\gamma_k$ are bounded; and $S_1(t)=\int_0^tS(u)\,du=O(\log t)$.
\end{theorem}

\begin{theorem}[Bellotti--Wong {\cite{BellottiWong2025}}]\label{thm:known_BW}
For $T\ge e$,
\begin{equation}
\Bigl|N(T)-\frac T{2\pi}\log\frac T{2\pi e}\Bigr|\le0.10076\log T+0.24460\log\log T+8.08344 .
\label{eq:BW_bound}
\end{equation}
\end{theorem}
Earlier explicit forms are due to Trudgian \cite{Trudgian2014} and Hasanalizade, Shen and Wong \cite{HasanalizadeShenWong2022}.

\begin{theorem}[Selberg {\cite{Selberg1946a}}; see {\cite[p.~225]{Fujii1974}} and {\cite[\S9.25]{Titchmarsh1986}}]\label{thm:known_selberg_moments}
For every $j\in\mathbb Z_{\ge1}$, with $c_j=\frac{(2j)!}{j!\,(2\pi)^{2j}}$,
\begin{equation}
\int_T^{2T}S(t)^{2j}\,dt=c_j\,T(\log\log T)^j+O_j\bigl(T(\log\log T)^{j-1/2}\bigr).
\label{eq:selberg_moments}
\end{equation}
\end{theorem}
In \cite[p.~225]{Fujii1974} the formula is quoted for $\int_T^{T+H}$ with $T^{1/2+a}\le H\le T$; we use $H=T$.

\begin{theorem}[Selberg's approximation, in the form of {\cite[\S3, pp.~232--233]{Fujii1974}}; \cite{Selberg1946b}]\label{thm:known_selberg_approx}
Let $x_0=T^{1/160}$ and $S^{\rm pr}_{x_0}(t)=-\frac1\pi\sum_{p<x_0}p^{-1/2}\sin(t\log p)$. Then
\begin{equation}
\int_T^{2T}\bigl(S(t)-S^{\rm pr}_{x_0}(t)\bigr)^4dt\ll T .
\label{eq:selberg_approx}
\end{equation}
\end{theorem}

\begin{proof}[Reference to the proof]
In \cite[\S3]{Fujii1974}, for a primitive character $\chi$ modulo $q$ and $x\ge2$, Selberg's expression $S(t,\chi)-\frac1\pi\Im\sum_{p<x^3}\chi(p)p^{-1/2-it}=\sum_{i=1}^5O(B_i(t))$ is written with explicit functions $B_1,\dots,B_5$, and it is shown on p.~233 that $\|B_i\|\ll k^2H^{1/2k}$ for $x=T^{a_2/60k}$, where $\|f\|=(\int_T^{T+H}|f|^{2k})^{1/2k}$ and the hypotheses are those of the Main Theorem of \cite{Fujii1974}: $0<a_1,a_2\le\frac14$, $q\le T^{1/4-a_1}$ and $T^{1/2+a_2}\le H\le T$. Take $\chi$ principal ($q=1$, $\zeta=L(s,\chi_0)$), $k=2$, $a_2=\frac14$ and $H=T$. Then $x^3=T^{1/160}$, and $\frac1\pi\Im p^{-1/2-it}=-\frac1\pi p^{-1/2}\sin(t\log p)$.
\end{proof}

\begin{theorem}[Montgomery--Vaughan {\cite[Cor.~3]{MontgomeryVaughan1974}}]\label{thm:known_MV}
For complex $b_n$, $\int_T^{2T}\bigl|\sum_{n\le Y}b_nn^{-it}\bigr|^2dt=\sum_{n\le Y}|b_n|^2\bigl(T+O(n)\bigr)$ with an absolute implied constant.
\end{theorem}
The version on $[T,2T]$ follows from the one on $[0,T]$ by replacing $b_n$ by $b_nn^{-iT}$.

\begin{theorem}[Mertens {\cite{Mertens1874}}]\label{thm:known_mertens}
For $y\ge2$, $\sum_{p\le y}\frac{\log p}p=\log y+O(1)$ and $\sum_{p\le y}\frac1p=\log\log y+B_1+O(\frac1{\log y})$ with a constant $B_1$.
\end{theorem}

\begin{theorem}[Kusmin--Landau; {\cite[Thm.~2.1]{GrahamKolesnik1991}}]\label{thm:known_KL}
Let $f$ be real with a continuous monotone derivative on $[a,b]$, and suppose that $\|f'(x)\|\ge\lambda>0$ there, where $\|\cdot\|$ is the distance to the nearest integer. Then $\sum_{a<n\le b}e^{2\pi if(n)}\ll\lambda^{-1}$.
\end{theorem}

\begin{theorem}[Landau {\cite{Landau1912}}; uniform version by Gonek {\cite{Gonek1993}}]\label{thm:known_landau}
For fixed $x>1$, $\sum_{0<\gamma\le T}x^{\rho}=-\frac T{2\pi}\Lambda(x)+O(\log T)$, where $\Lambda(x)=0$ unless $x$ is a prime power, and the zeros are counted with multiplicity.
\end{theorem}

\begin{theorem}[Littlewood; {\cite[Thm.~14.13]{Titchmarsh1986}}]\label{thm:known_RH_S}
On the Riemann hypothesis, $S(t)=O(\log t/\log\log t)$.
\end{theorem}

\section{Identities at the zero ordinates}
\label{sec:identities}

The identities below are exact and involve no hypothesis on the zeros, not even their simplicity. Their specific application to ordinal defects is the starting point of the paper; the broader quadrature and moment methods are classical.

\begin{theorem}\label{thm:identity_general}
Let $m,n\in\mathbb Z_{\ge1}$ with $m<n$. Let $F:[\gamma_m,\gamma_n]\to\mathbb R$ be of class $C^1$ on $[\gamma_m,\gamma_n]$ and put $u_k=k-\frac12-F(\gamma_k)$ for $m\le k\le n$. For every integer $j\ge1$,
\begin{equation}
\int_{\gamma_m}^{\gamma_n}\bigl(N(t)-F(t)\bigr)^{2j}F'(t)\,dt=\sum_{k=m+1}^{n-1}P_j(u_k)+\frac{(\frac12+u_m)^{2j+1}+(\frac12-u_n)^{2j+1}}{2j+1},
\label{eq:identity_general}
\end{equation}
where
\begin{equation}
P_j(u)=\frac{(\frac12+u)^{2j+1}+(\frac12-u)^{2j+1}}{2j+1}=\sum_{i=0}^{j}\binom{2j+1}{2i}\frac{2^{2i-2j}}{2j+1}\,u^{2i}=u^{2j}+\sum_{i<j}b_{j,i}u^{2i}.
\label{eq:Pj_def}
\end{equation}
In particular, for $j=1$,
\begin{equation}
\sum_{k=m+1}^{n-1}u_k^2=\int_{\gamma_m}^{\gamma_n}\bigl(N(t)-F(t)\bigr)^2F'(t)\,dt-\frac{n-m-1}{12}-\frac{(\frac12+u_m)^3+(\frac12-u_n)^3}3 .
\label{eq:identity_quadratic}
\end{equation}
\end{theorem}

\begin{proof}
Fix $k$ with $m\le k<n$. If $\gamma_k<\gamma_{k+1}$ then exactly $k$ ordinates, counted with multiplicity, are $\le t$ for $t\in(\gamma_k,\gamma_{k+1})$, so $N(t)=k$ there. Define $H_k(y)=-(k-y)^{2j+1}/(2j+1)$.  Since $F\in C^1$, the chain rule and the fundamental theorem of calculus give
\begin{equation}
\int_{\gamma_k}^{\gamma_{k+1}}\bigl(k-F(t)\bigr)^{2j}F'(t)\,dt=\frac{(\frac12-u_{k+1})^{2j+1}+(\frac12+u_k)^{2j+1}}{2j+1},
\label{eq:identity_one_gap}
\end{equation}
because $F(\gamma_k)-k=-\frac12-u_k$ and $F(\gamma_{k+1})-k=\frac12-u_{k+1}$. If $\gamma_k=\gamma_{k+1}$ the left side of \eqref{eq:identity_one_gap} is $0$; since then $u_{k+1}=u_k+1$, also the right side is $(-\frac12-u_k)^{2j+1}+(\frac12+u_k)^{2j+1}=0$. The values of $N$ at the finitely many ordinates do not affect the integrals. Summing \eqref{eq:identity_one_gap} over $m\le k<n$, the term $(\frac12+u_k)^{2j+1}$ with $m<k<n$ meets the term $(\frac12-u_k)^{2j+1}$ coming from the gap $(\gamma_{k-1},\gamma_k)$, which gives \eqref{eq:identity_general}. The expansion \eqref{eq:Pj_def} keeps the even powers of the binomial expansion, and the coefficient of $u^{2j}$ is $\binom{2j+1}{2j}/(2j+1)=1$. For $j=1$, $P_1(u)=u^2+\frac1{12}$.
\end{proof}

\begin{remark}\label{rem:twelfth}
The constant $\frac1{12}$ is the mean square of a centred sawtooth of unit step: between two consecutive zeros $N-F$ decreases through one unit in the variable $y=F(t)$. Identity \eqref{eq:identity_quadratic} therefore says that the mean square of the defect at the zeros is the $dF$-mean square of $N-F$ minus the contribution of the sawtooth. Nothing is assumed about the sign of $F'$; when $F'$ changes sign the identity remains exact, with the integral counted with sign.
\end{remark}

\begin{corollary}\label{cor:identity_clock}
With $F=\NO$, so that $u_k=s_k$ and $N-\NO=S+\omega/\pi$ almost everywhere,
\begin{equation}
\sum_{k=m+1}^{n-1}s_k^2=\int_{\gamma_m}^{\gamma_n}\Bigl(S(t)+\frac{\omega(t)}\pi\Bigr)^2\lL(t)\,dt-\frac{n-m-1}{12}-\frac{(\frac12+s_m)^3+(\frac12-s_n)^3}3 .
\label{eq:identity_clock}
\end{equation}
With $F=F_X$, so that $u_k=u_k(X)$ and $N-F_X=S-S_X+\omega/\pi$ almost everywhere,
\begin{align}
\sum_{k=m+1}^{n-1}u_k(X)^2
&=\int_{\gamma_m}^{\gamma_n}\Bigl(S(t)-S_X(t)+\frac{\omega(t)}\pi\Bigr)^2
  \bigl(\lL(t)+S_X'(t)\bigr)\,dt \\
&\quad-\frac{n-m-1}{12}
-\frac{(\frac12+u_m)^3+(\frac12-u_n)^3}{3}.
\label{eq:identity_LeClair}
\end{align}
\end{corollary}

\begin{proof}
Theorem~\ref{thm:identity_general} with $j=1$ and \eqref{eq:N_decomposition}, \eqref{eq:defect_def}.
\end{proof}

\begin{lemma}\label{lem:first_moment_identity}
Let $2<a<b$ be points which are not ordinates, let $f\in C^1[a,b]$ and $G=N-\NO$. Then
\begin{equation}
\sum_{a<\gamma_k\le b}s_kf(\gamma_k)=\int_a^bG(t)f(t)\lL(t)\,dt+\frac12\Bigl[f\,G^2\Bigr]_a^b-\frac12\int_a^bG(t)^2f'(t)\,dt .
\label{eq:first_moment_identity}
\end{equation}
\end{lemma}

\begin{proof}
Let $\gamma$ be an ordinate of multiplicity $\nu$ and $G^*(\gamma)=\frac12(G(\gamma^-)+G(\gamma^+))$. The indices with $\gamma_k=\gamma$ are $N(\gamma^-)+1,\dots,N(\gamma^-)+\nu$, so $\sum_{\gamma_k=\gamma}(k-\frac12)=\nu N(\gamma^-)+\frac{\nu^2}2=\nu\,\frac{N(\gamma^-)+N(\gamma^+)}2$ and $\sum_{\gamma_k=\gamma}s_k=\nu\,G^*(\gamma)$. Hence the left side of \eqref{eq:first_moment_identity} is the Stieltjes integral $\int_{(a,b]}G^*f\,dN$. As measures $dN=dG+\lL\,dt$. The function $G$ is $C^1$ between ordinates and jumps by $\nu$ at an ordinate of multiplicity $\nu$; at a jump $\Delta(G^2)=2G^*\Delta G$, and between jumps $(G^2)'=2GG'$. Hence $G^*\,dG=\frac12d(G^2)$, and integration by parts gives $\int_{(a,b]}f\,d(G^2)=[fG^2]_a^b-\int_a^bG^2f'\,dt$.
\end{proof}

\begin{corollary}\label{cor:cumulative}
$\sum_{\gamma_k\le t}s_k=O(\log^2t)$. In particular the mean of $s_k$ over $k\le K$ is $O(K^{-1}\log^2K)$.
\end{corollary}

\begin{proof}
Apply Lemma~\ref{lem:first_moment_identity} with $f\equiv1$, $a=14$ (below $\gamma_1$) and $b=t$. Then $G^2=O(\log^2t)$ by Theorem~\ref{thm:known_basic}, and $\int_{14}^tG\lL\,du=\int_{14}^tS\lL\,du+O(\log^2t)=\lL(t)S_1(t)-\int_{14}^tS_1(u)\frac{du}{2\pi u}+O(\log t)=O(\log^2t)$ by $S_1=O(\log t)$.
\end{proof}

\section{Moments and extremes of the clock displacement}
\label{sec:clock_moments}

\begin{theorem}\label{thm:clock_moments}
For every fixed integer $j\ge1$, as $T\to\infty$,
\begin{equation}
\frac1{\#\IT}\sum_{k\in\IT}s_k^{2j}=c_j(\log\log T)^j+O_j\bigl((\log\log T)^{j-1/2}\bigr),\qquad c_j=\frac{(2j)!}{j!\,(2\pi)^{2j}} .
\label{eq:clock_moments}
\end{equation}
In particular the mean square of $s_k$ over $\IT$ is $\frac1{2\pi^2}\log\log T+O(\sqrt{\log\log T})$ and
\begin{equation}
\frac1{\#\IT}\sum_{k\in\IT}\delta_k^2=\frac{2\log\log T}{\log^2T}\bigl(1+O((\log\log T)^{-1/2})\bigr).
\label{eq:delta_mean_square}
\end{equation}
\end{theorem}

\begin{proof}
Let $m$ and $n$ be the smallest and largest indices in $\IT$. Theorem~\ref{thm:identity_general} with $F=\NO$ gives
\begin{equation}
\sum_{k=m+1}^{n-1}P_j(s_k)=\int_{\gamma_m}^{\gamma_n}\Bigl(S+\frac\omega\pi\Bigr)^{2j}\lL\,dt+O(\log^{2j+1}T),
\label{eq:moments_step1}
\end{equation}
since $|s_k|\ll\log T$ by Lemma~\ref{lem:s_bounds}. The intervals $[T,\gamma_m]$ and $[\gamma_n,2T]$ have bounded length (Theorem~\ref{thm:known_basic}) and the integrand is $O(\log^{2j+1}T)$ there; moreover $|(S+\omega/\pi)^{2j}-S^{2j}|\ll T^{-1}\log^{2j}T$. Since $\lL(t)=\lL(T)+O(1)$ on $[T,2T]$, Theorem~\ref{thm:known_selberg_moments} yields
\begin{equation}
\int_T^{2T}S^{2j}\lL\,dt=\lL(T)\,c_jT(\log\log T)^j+O\bigl(\lL(T)\,T(\log\log T)^{j-1/2}\bigr).
\label{eq:moments_step2}
\end{equation}
Adding the terms $k=m,n$, which are $O(\log^{2j}T)$, and dividing by \eqref{eq:IT_size},
\begin{equation}
\frac1{\#\IT}\sum_{k\in\IT}P_j(s_k)=c_j(\log\log T)^j+O\bigl((\log\log T)^{j-1/2}\bigr).
\label{eq:moments_step3}
\end{equation}
For $j=1$ this is \eqref{eq:clock_moments}, because $P_1(s)=s^2+\frac1{12}$. If \eqref{eq:clock_moments} holds for all $i<j$, then $\sum_{i<j}b_{j,i}\frac1{\#\IT}\sum s_k^{2i}=O((\log\log T)^{j-1})$ and \eqref{eq:clock_moments} follows for $j$. For \eqref{eq:delta_mean_square} use \eqref{eq:s_delta}: $\lL(\xi_k)=\lL(T)(1+O(1/\log T))$ because $|\delta_k|\ll1$, and $\lL(T)^2=(2\pi)^{-2}\log^2T\,(1+O(1/\log T))$.
\end{proof}

\begin{corollary}\label{cor:clock_distribution}
Let $\sigma_T^2=\frac1{2\pi^2}\log\log T$ and let $Z$ be a standard normal variable. For every fixed $\eta\ge0$,
\begin{equation}
\frac1{\#\IT}\#\{k\in\IT:|s_k|\le\eta\,\sigma_T\}\longrightarrow\mathbb P(|Z|\le\eta).
\label{eq:clock_distribution}
\end{equation}
In particular, for every fixed $A>0$, $\#\{k\in\IT:|s_k|\le A\}=o(\#\IT)$.
\end{corollary}

\begin{proof}
By Theorem~\ref{thm:clock_moments} the empirical distributions of $s_k^2/\sigma_T^2$, $k\in\IT$, have $j$-th moments tending to $c_j(2\pi^2)^j=\frac{(2j)!}{j!\,2^j}=\mathbb E[Z^{2j}]$. The law of $Z^2$ is determined by its moments, its moment generating function being finite near $0$, so the method of moments \cite[\S30]{Billingsley1995} gives convergence in distribution of $s_k^2/\sigma_T^2$ to $Z^2$; the limit law is continuous.
\end{proof}

\begin{remark}\label{rem:korolev}
Selberg stated, and Korolev proved \cite{Korolev2012}, moment formulae for the integer displacement of $\gamma_n$ with respect to the Gram intervals, and Korolev deduced that $\gamma_n$ lies in its own Gram interval only for a set of indices of density zero. By Remark~\ref{rem:gram_points}, Theorem~\ref{thm:clock_moments} and Corollary~\ref{cor:clock_distribution} are the analogous statements for the real displacement $s_k$ with respect to the half-integer Gram points $\tau_k$. The proof given here is a direct consequence of Theorem~\ref{thm:identity_general} and Selberg's moments.
\end{remark}

\begin{proposition}\label{prop:extremes}
\begin{enumerate}[label=(\alph*),leftmargin=2em]
\item $\limsup_{k\to\infty}s_k=+\infty$ and $\liminf_{k\to\infty}s_k=-\infty$.
\item $\limsup_{k\to\infty}|\delta_k|\le2\pi\cdot0.10076=0.633093\ldots$.
\item On the Riemann hypothesis, $\delta_k\to0$.
\end{enumerate}
\end{proposition}

\begin{proof}
(a) Let $\gamma_k<\gamma_{k+1}$ and $t\in(\gamma_k,\gamma_{k+1})$. Then $N(t)=k$ and, by monotonicity of $\NO$,
\begin{equation}
s_{k+1}-\frac12\le N(t)-\NO(t)\le s_k+\frac12 .
\label{eq:sawtooth_sandwich}
\end{equation}
Suppose $s_k\le A$ for all large $k$. Then $S(t)\le A+1$ for all large $t$ by \eqref{eq:N_decomposition}. By Theorem~\ref{thm:known_basic}, $\int_T^{2T}S\,dt=O(\log T)$, so $\int_T^{2T}S^-dt=\int_T^{2T}S^+dt-\int_T^{2T}S\,dt\le(A+1)T+O(\log T)$ and $\int_T^{2T}|S|\,dt\ll T$. Hölder's inequality with exponents $\frac32$ and $3$ and Theorem~\ref{thm:known_selberg_moments} with $j=2$ give
\begin{equation}
\int_T^{2T}S^2dt\le\Bigl(\int_T^{2T}|S|\,dt\Bigr)^{2/3}\Bigl(\int_T^{2T}S^4dt\Bigr)^{1/3}\ll T(\log\log T)^{2/3},
\label{eq:holder_contradiction}
\end{equation}
which contradicts \eqref{eq:selberg_moments} with $j=1$. The statement for $\liminf$ follows in the same way from the left inequality in \eqref{eq:sawtooth_sandwich}.

(b) By Lemma~\ref{lem:s_bounds}, $|s_k|\le\max_\pm|N(\gamma_k^\pm)-\NO(\gamma_k)|+\frac12$, and $\NO(T)=\frac T{2\pi}\log\frac T{2\pi e}+\frac78$. Theorem~\ref{thm:known_BW} gives $|s_k|\le0.10076\log\gamma_k+O(\log\log\gamma_k)$. By \eqref{eq:s_delta}, $|\delta_k|=|s_k|/\lL(\xi_k)$ with $\lL(\xi_k)=\frac1{2\pi}\log\gamma_k\,(1+o(1))$.

(c) Theorem~\ref{thm:known_RH_S} and Lemma~\ref{lem:s_bounds} give $s_k=O(\log\gamma_k/\log\log\gamma_k)$, hence $\delta_k=O(1/\log\log\gamma_k)$.
\end{proof}

\begin{proposition}[Effective finite-height consequence of Bellotti--Wong]\label{prop:BW_pointwise}
Let $E(t)=0.10076\log t+0.24460\log\log t+8.08344$. For every $k$,
\begin{equation}
 -\frac{E(\gamma_k)+3/8}{\lL(\tau_k)}
 \le\delta_k\le\frac{E(\gamma_k)-11/8}{\lL(\gamma_k)}.
 \label{eq:BW_pointwise}
\end{equation}
These inequalities are analytic and unconditional, and count repeated ordinates individually.
\end{proposition}
\begin{proof}
Apply \eqref{eq:BW_bound} to the one-sided counts at $\gamma_k$ and use Lemma~\ref{lem:s_bounds}. The smooth count exceeds the central counting term in that theorem by $7/8$, so one gets
\begin{equation}
 -E(\gamma_k)-3/8\le s_k\le E(\gamma_k)-11/8.
 \label{eq:BW_ordinal_sandwich}
\end{equation}
When $\delta_k\ge0$, monotonicity of $\lL$ and $s_k=\int_{\gamma_k}^{\tau_k}\lL(t)\,dt$ give the upper bound in \eqref{eq:BW_pointwise}. When $\delta_k\le0$, $-s_k=\int_{\tau_k}^{\gamma_k}\lL(t)\,dt\ge-\delta_k\lL(\tau_k)$ gives the lower bound. The other inequality in each case is automatic, as $E(\gamma_k)>11/8$ and $\tau_k,\gamma_k>2\pi$.
\end{proof}
The bound of Bellotti--Wong for the counting error grows with $t$. Dividing after inversion by the local density of zeros gives the finite ordinate limit superior $2\pi\cdot0.10076=0.6330937515\ldots$. This consequence is deduced from their theorem; it is not a stronger independent estimate for their counting function. A limit superior is also distinct from $\sup_{k\ge1}|\delta_k|$: the measured value $|\delta_{34}|=0.993711$ already exceeds $0.6331$. Proposition~\ref{prop:BW_pointwise} is a finite-height envelope, not an assertion that $0.6331$ bounds every ordinate displacement.

\begin{proposition}\label{prop:fixed_sup}
Let $X\ge2$ be fixed and $B\in(0,\infty]$. Then $\limsup_k|\delta_k^{X,B}|=\limsup_k|\delta_k|$, while $\limsup_k|u_k(X)|=\infty$.
\end{proposition}

\begin{proof}
$|\delta_k^{X,B}-\delta_k|\le|S_X(\tau_k)|/\lL(\tau_k)\le\sigma_X/\lL(\tau_k)\to0$ by \eqref{eq:sigma_M_def}, and $|u_k(X)-s_k|\le\sigma_X$; apply Proposition~\ref{prop:extremes}(a).
\end{proof}

\section{The main theorem: a variable cutoff in the truncated equation}
\label{sec:main_theorem}

\subsection{Statement}
\label{subsec:main_statement}

\begin{theorem}\label{thm:main}
There exist real absolute constants $T_0>2\pi$ and $C>0$ such that, for every real $T\ge T_0$ and every real $X\in[2,T^{1/2}]$, writing $c=\log X/\log T$,
\begin{equation}
\Bigl|\frac1{\#\IT}\sum_{k\in\IT}u_k(X)^2-\frac1{2\pi^2}\log\frac1c\Bigr|\le C\Bigl(1+\sqrt{\log(1/c)}\Bigr).
\label{eq:main_theorem}
\end{equation}
\end{theorem}

The proof is given in Subsection~\ref{subsec:main_proof}. For fixed $c>0$, both the displayed main term and its error are bounded quantities: the estimate does not identify a limiting constant or an asymptotic equivalence in that regime. It controls the defect at the true ordinate, and does not select a root of $F_X(t)=k-\frac12$ when $F_X$ is not monotone. Note that $\log(1/c)=\log\log T-\log\log X$.

\begin{corollary}[fixed cutoff]\label{cor:fixed_cutoff}
Let $X\ge2$ be fixed. Then
\begin{equation}
\frac1{\#\IT}\sum_{k\in\IT}u_k(X)^2=\frac1{2\pi^2}\log\log T+O_X\bigl(\sqrt{\log\log T}\bigr),
\label{eq:fixed_cutoff}
\end{equation}
and for every fixed $A>0$, $\#\{k\in\IT:|u_k(X)|\le A\}=o(\#\IT)$.
\end{corollary}

\begin{proof}
The first statement is \eqref{eq:main_theorem}. For the second, $|u_k(X)-s_k|\le\sigma_X$, so $|u_k(X)|\le A$ implies $|s_k|\le A+\sigma_X$; apply Corollary~\ref{cor:clock_distribution}.
\end{proof}

\begin{corollary}[power cutoff]\label{cor:power_cutoff}
Let $0<c\le\frac12$ be fixed and $X=T^c$. Then
\begin{equation}
\limsup_{T\to\infty}\frac1{\#\IT}\sum_{k\in\IT}u_k(T^c)^2\le\frac1{2\pi^2}\log\frac1c+C\Bigl(1+\sqrt{\log(1/c)}\Bigr),
\label{eq:power_cutoff}
\end{equation}
and for every function $\eta(T)\to\infty$, $\#\{k\in\IT:|u_k(T^c)|>\eta(T)\}=o(\#\IT)$. The same holds for $X=K^c$ with $K=N(T)$ and fixed $0<c<\frac12$, because then $X=T^{c_T}$ with $c_T=c+O(\log\log T/\log T)$, so that $c_T<\frac12$ for large $T$ and $\log(1/c_T)=\log(1/c)+o(1)$.
\end{corollary}

\begin{proof}
The bound is \eqref{eq:main_theorem}; the second statement follows by Chebyshev's inequality.
\end{proof}

\begin{corollary}[dichotomy]\label{cor:dichotomy}
Let $X=X(T)$ with $2\le X\le T^{1/2}$. The mean square of $u_k(X)$ over $k\in\IT$ is bounded as $T\to\infty$ if and only if $\liminf_{T\to\infty}\log X/\log T>0$. In particular, for any fixed $X_0\ge2$ and fixed $c\in(0,\frac12]$,
\begin{equation}
\frac1{\#\IT}\sum_{k\in\IT}\Bigl(u_k(X_0)^2-u_k(T^c)^2\Bigr)=\frac1{2\pi^2}\log\log T+O\bigl(\sqrt{\log\log T}\bigr)\longrightarrow\infty .
\label{eq:dichotomy}
\end{equation}
\end{corollary}

\begin{proof}
Both directions follow from \eqref{eq:main_theorem}, since $x\mapsto\frac{x}{2\pi^2}-C(1+\sqrt x)$ tends to infinity. For \eqref{eq:dichotomy} subtract \eqref{eq:power_cutoff} from \eqref{eq:fixed_cutoff}.
\end{proof}

\begin{remark}\label{rem:comparison_LeClair}
Corollary~\ref{cor:dichotomy} is the precise sense in which an ordinal-dependent number of primes improves the truncated zero-counting equation. With $N=\#\{p\le X\}$ primes, $\log N=\log X+O(\log\log X)$, so the condition is that $\log N$ grows proportionally to $\log T$. A fixed $N$, as in the computations of \cite[Sect.~VI]{LeClair2021}, leaves a defect whose mean square grows like $(2\pi^2)^{-1}\log\log T$, and almost all zeros then violate the truncated equation by more than any fixed amount. LeClair's requirement $N(t)\to\infty$ \cite[Sect.~III]{LeClair2021} is necessary but not sufficient: if $\log N=o(\log T)$, the mean-square defect still tends to infinity. The range $c>\frac12$, which contains LeClair's upper limit $N<[t^2]$, is not covered by Theorem~\ref{thm:main}.
\end{remark}

\begin{remark}\label{rem:block_rule}
Theorem~\ref{thm:main} concerns a cutoff that is constant on each window $\IT$. Taking $T=2^r$ and $X=2^{rc}$ on $\mathcal I_{2^r}$ gives a cutoff depending on the ordinal through the window; the rule $X_k=(\tau_k/2\pi)^\theta$ used in Section~\ref{sec:selected_computations} differs from such a block rule by a factor at most $2^\theta$ inside each window. The extension of \eqref{eq:main_theorem} to cutoffs that vary from zero to zero is not proved here.
\end{remark}

\subsection{Lemmas on prime polynomials}
\label{subsec:main_lemmas}

Fix $T$ large and $2\le X\le T^{1/2}$, let $x_0=T^{1/160}$ and decompose $S_X=S_X^{(1)}+S_X^{(2)}+S_X^{(\ge3)}$ according to $m=1$, $m=2$ and $m\ge3$ in \eqref{eq:SX_def}. Let $\mathcal P$ be the set of primes in $(X,x_0)$ if $X<x_0$, and in $[x_0,X]$ if $X\ge x_0$, let $\epsilon=+1$ in the first case and $\epsilon=-1$ in the second, and put
\begin{equation}
D_1(t)=-\frac1\pi\sum_{p\in\mathcal P}\frac{\sin(t\log p)}{\sqrt p},\qquad V=\sum_{p\in\mathcal P}\frac1p,\qquad D=S^{\rm pr}_{x_0}-S_X=\epsilon D_1-S_X^{(2)}-S_X^{(\ge3)} .
\label{eq:D_def}
\end{equation}

\begin{lemma}\label{lem:V_size}
$V=\log(1/c)+O(1)$ uniformly for $0<c\le\frac12$ with $X=T^c\ge2$.
\end{lemma}

\begin{proof}
By Theorem~\ref{thm:known_mertens}, $V=|\log\log x_0-\log\log X|+O(1)=|\log\frac{1}{160c}|+O(1)$. If $c\le\frac1{160}$ this is $\log\frac1c-\log160+O(1)$; if $\frac1{160}<c\le\frac12$ both $V$ and $\log\frac1c$ are $O(1)$.
\end{proof}

\begin{lemma}\label{lem:trig_MV}
For real $b_n$, $2\le n\le Y$,
\begin{equation}
\int_T^{2T}\Bigl(\sum_{n\le Y}b_n\sin(t\log n)\Bigr)^2dt=\frac T2\sum_{n\le Y}b_n^2+O\Bigl(\sum_{n\le Y}nb_n^2+\Bigl(\sum_{n\le Y}|b_n|\Bigr)^2\Bigr),
\label{eq:trig_MV}
\end{equation}
and the same holds with $\cos$ in place of $\sin$.
\end{lemma}

\begin{proof}
With $Q(t)=\sum b_nn^{-it}$ we have $\sum b_n\sin(t\log n)=-\Im Q$ and $(\Im Q)^2=\frac12(|Q|^2-\Re Q^2)$. Theorem~\ref{thm:known_MV} evaluates $\int|Q|^2$, and $|\int_T^{2T}(mn)^{-it}dt|\le2/\log(mn)\le2/\log4$ bounds $\int Q^2$ by $O((\sum|b_n|)^2)$. For $\cos$ use $\Re Q$.
\end{proof}

\begin{lemma}\label{lem:D_moments}
With absolute implied constants,
\begin{align}
\int_T^{2T}D^2\,dt&=\frac T{2\pi^2}V+O\bigl(T(1+\sqrt V)\bigr),
\label{eq:D_second}\\
\int_T^{2T}D^4\,dt&\ll T(1+V)^2,
\label{eq:D_fourth}\\
\int_T^{2T}S_X'(t)^2\,dt&\ll T(1+\log X)^2 .
\label{eq:SX_derivative}
\end{align}
\end{lemma}

\begin{proof}
All primes involved are $\le\max(X,x_0)\le T^{1/2}$. By Lemma~\ref{lem:trig_MV} with $b_p=\frac1{\pi\sqrt p}$, $\int D_1^2=\frac T{2\pi^2}V+O(T^{1/2}+(\sum_{p\le T^{1/2}}p^{-1/2})^2)=\frac T{2\pi^2}V+O(T^{1/2})$. For $S_X^{(2)}(t)=-\frac1{2\pi}\sum_{p\le X}p^{-1}\sin(t\log p^2)$, Lemma~\ref{lem:trig_MV} gives $\int(S_X^{(2)})^2\ll T\sum p^{-2}+\#\{p\le X\}+(\sum_{p\le X}p^{-1})^2\ll T$. Finally $|S_X^{(\ge3)}|\le\frac1\pi\sum_p\sum_{m\ge3}\frac1{mp^{m/2}}<\infty$. Expanding $D^2$ and using the Cauchy--Schwarz inequality for the cross terms gives \eqref{eq:D_second}.

For \eqref{eq:D_fourth}, $(a+b+c)^4\le27(a^4+b^4+c^4)$ reduces it to the three parts. Let $Q(t)=\sum_{p\in\mathcal P}p^{-1/2}p^{-it}$, so $|D_1|\le|Q|/\pi$ and $Q^2=\sum_\nu c_\nu \nu^{-it}$ with $c_{p^2}=p^{-1}$, $c_{pq}=2(pq)^{-1/2}$ for $p<q$, and $\nu\le T$. Theorem~\ref{thm:known_MV} gives $\int|Q|^4=\sum|c_\nu|^2(T+O(\nu))\ll T\sum|c_\nu|^2\le T(1+2V^2)$. For $S_X^{(2)}$ the same argument applies to $Q_2(t)=\sum_{p\le X}p^{-1}(p^2)^{-it}$, since $|S_X^{(2)}|\le|Q_2|/2\pi$: $Q_2^2=\sum d_\nu\nu^{-it}$ with $\nu=(pq)^2$, $|d_\nu|\le2(pq)^{-1}$, so $\sum|d_\nu|^2\ll1$ and $\sum\nu|d_\nu|^2\le4\#\{p\le X\}^2\le4X^2\le4T$. The part $S_X^{(\ge3)}$ is bounded.

For \eqref{eq:SX_derivative}, $S_X'(t)=\sum_{n\in A_X}a_n\log n\cos(t\log n)$. The terms $m=1$ have coefficients $\frac{\log p}{\pi\sqrt p}$, and Lemma~\ref{lem:trig_MV} with Theorem~\ref{thm:known_mertens} gives $\ll T\log X\sum_{p\le X}\frac{\log p}p+\sum_{p\le X}\log^2p+(\sum_{p\le X}\frac{\log p}{\sqrt p})^2\ll T(1+\log X)^2+X\log^2X$, and $X\le T^{1/2}$. The terms $m=2$ have coefficients $\frac{\log p}{\pi p}$ and contribute $\ll T+X\log X+\log^2X$. The terms $m\ge3$ are bounded by $\frac1\pi\sum_p\sum_{m\ge3}p^{-m/2}\log p<\infty$.
\end{proof}

\subsection{Proof of the main theorem}
\label{subsec:main_proof}

Let $E_0=S-S^{\rm pr}_{x_0}$ and $R=S-S_X=E_0+D$. By Theorem~\ref{thm:known_selberg_approx} and Hölder's inequality,
\begin{equation}
\int_T^{2T}E_0^4\,dt\ll T,\qquad\int_T^{2T}E_0^2\,dt\le T^{1/2}\Bigl(\int_T^{2T}E_0^4\,dt\Bigr)^{1/2}\ll T .
\label{eq:E0_bounds}
\end{equation}
Write $J=\int_T^{2T}R^2\,dt$. From \eqref{eq:D_second}, \eqref{eq:E0_bounds} and the Cauchy--Schwarz inequality,
\begin{equation}
J=\frac T{2\pi^2}V+O\bigl(T(1+\sqrt V)\bigr),
\label{eq:J_asymptotic}
\end{equation}
and from \eqref{eq:D_fourth}, \eqref{eq:E0_bounds} and $(a+b)^4\le8(a^4+b^4)$,
\begin{equation}
\int_T^{2T}R^4\,dt\ll T(1+V)^2 .
\label{eq:R_fourth}
\end{equation}

Let $m$ and $n$ be the smallest and largest indices in $\IT$ and apply \eqref{eq:identity_LeClair}. By Lemma~\ref{lem:s_bounds} and \eqref{eq:sigma_M_def}, $|u_k(X)|\ll\log T+\sqrt X\ll T^{1/4}$, so the boundary terms and $u_m^2,u_n^2$ are $O(T^{3/4})$. The intervals $[T,\gamma_m]$ and $[\gamma_n,2T]$ have bounded length and on them $|R|^2(\lL+|S_X'|)\ll T^{3/4}$. Since $|\omega|\ll1/T$, replacing $(R+\omega/\pi)^2$ by $R^2$ costs $\ll T^{-1}\int_T^{2T}(|R|+1)(\lL+|S_X'|)\,dt\ll(1+V)^{1/2}\log T$ by the Cauchy--Schwarz inequality, \eqref{eq:J_asymptotic} and \eqref{eq:SX_derivative}. Hence
\begin{equation}
\sum_{k\in\IT}u_k(X)^2=\int_T^{2T}R^2\lL\,dt+\int_T^{2T}R^2S_X'\,dt-\frac{\#\IT}{12}+O(T^{3/4}).
\label{eq:main_step1}
\end{equation}
Since $0\le\lL(t)-\lL(T)\le\frac{\log2}{2\pi}$ on $[T,2T]$, the first integral is $\lL(T)J+O(J)$. By the Cauchy--Schwarz inequality, \eqref{eq:R_fourth} and \eqref{eq:SX_derivative},
\begin{equation}
\Bigl|\int_T^{2T}R^2S_X'\,dt\Bigr|\ll T(1+V)(1+\log X).
\label{eq:main_step2}
\end{equation}
After division by $\#\IT\asymp T\log T$, this contributes
\[
 O\!\left(\frac{1+V}{\log T}+c(1+V)\right).
\]
Here $0<c\le1/2$, $V=\log(1/c)+O(1)$, and $c(1+\log(1/c))$ is bounded on $(0,1/2]$.  Moreover, $X\ge2$ implies $c\ge(\log2)/(\log T)$ and hence $V\ll1+\log\log T$, so $(1+V)/\log T=o(1)$.  Thus no unproved improvement of \eqref{eq:main_step2} is required.  Dividing by $\#\IT=T\lL(T)+O(T)$ and using \eqref{eq:J_asymptotic},
\begin{equation}
\frac1{\#\IT}\sum_{k\in\IT}u_k(X)^2=\frac V{2\pi^2}+O\Bigl(1+\sqrt V+\frac{1+V}{\log T}+(1+V)\,c\Bigr),
\label{eq:main_step3}
\end{equation}
because $\log X=c\log T$. By Lemma~\ref{lem:V_size}, $(1+V)c\ll c(1+\log\frac1c)\ll1$ for $0<c\le\frac12$, and $V\ll\log\log T$ because $c\ge\log2/\log T$, so $(1+V)/\log T=o(1)$. A second application of Lemma~\ref{lem:V_size} replaces $\frac V{2\pi^2}$ by $\frac1{2\pi^2}\log\frac1c$ and $\sqrt V$ by $\sqrt{\log(1/c)}+O(1)$. This proves \eqref{eq:main_theorem}.\qed

\begin{remark}\label{rem:alternative_route}
An alternative route to an upper bound of the form $\sum_{k\in\IT}u_k(T^c)^2\ll_c\#\IT$ averages a pointwise inequality for $|u_k|$ over windows $[\gamma_k,\gamma_k+h]$ with $h=1/\log T$, and controls the overlap of the windows through the multiplicity $M_h(t)=\#\{k:\gamma_k\le t\le\gamma_k+h\}\le1+N(t)-N(t-h)$, whose moments are bounded by the corollary to the Main Theorem of \cite{Fujii1974}. Such an argument needs fourth moments of $S-S_X$ and of $S_X'$ and gives only an upper bound. The identity of Theorem~\ref{thm:identity_general} makes the window multiplicities unnecessary, needs only the second moment of $S_X'$, and gives the uniform main-term estimate \eqref{eq:main_theorem}. In the parametrization $X=K^c$ with $K\asymp T\log T$ the condition $X^2\le T$ of the Montgomery--Vaughan step requires $c<\frac12$; in the parametrization $X=T^c$ the endpoint $c=\frac12$ is included.
\end{remark}

\subsection{Localization of the zeros for small cutoffs}
\label{subsec:localization}

\begin{proposition}\label{prop:localization}
Let $X\ge2$ with $M_X\le\frac12\lL(T)$. Then $F_X$ is strictly increasing on $[T,2T]$, for every $k\in\IT$ with $F_X(T)\le k-\frac12\le F_X(2T)$ the truncated equation $F_X(t)=k-\frac12$ has exactly one solution $t_k^X\in[T,2T]$, and
\begin{equation}
|\gamma_k-t_k^X|\le\frac{2\,|u_k(X)|}{\lL(T)} .
\label{eq:localization}
\end{equation}
\end{proposition}

\begin{proof}
By \eqref{eq:sigma_M_def}, $F_X'\ge\lL(T)-M_X\ge\frac12\lL(T)$ on $[T,2T]$, and $u_k(X)=F_X(t_k^X)-F_X(\gamma_k)$; apply the mean value theorem.
\end{proof}

\begin{remark}\label{rem:localization_scope}
Since $M_X\asymp\sqrt X$, Proposition~\ref{prop:localization} applies only for $X\ll\log^2T$, where $\log(1/c)\sim\log\log T$; it then turns \eqref{eq:main_theorem} into a mean-square localization of $\gamma_k$ around the solutions of the truncated counting equation, but with a growing constant. For $X=T^c$ the bound $M_X\le\frac12\lL(T)$ fails, $F_X$ need not be monotone and the truncated equation may have several solutions near $\gamma_k$; the defect $u_k(X)$ measures the quality of the equation without selecting a solution.
\end{remark}

\section{Results conditional on a Landau-type hypothesis}
\label{sec:conditional}

\begin{hypothesis}[L]\label{hyp:L}
For every fixed integer $n\ge2$, as $t\to\infty$,
\begin{equation}
\sum_{0<\gamma\le t}n^{i\gamma}=-\frac t{2\pi}\frac{\Lambda(n)}{\sqrt n}+o(t),
\label{eq:hyp_L}
\end{equation}
the zeros being counted with multiplicity.
\end{hypothesis}

\begin{proposition}\label{prop:RH_implies_L}
The Riemann hypothesis implies Hypothesis~\ref{hyp:L}, with error $O(\log t)$.
\end{proposition}

\begin{proof}
On the Riemann hypothesis $n^\rho=\sqrt n\,n^{i\gamma}$; apply Theorem~\ref{thm:known_landau}.
\end{proof}

\begin{remark}\label{rem:L_unconditional}
Unconditionally, Theorem~\ref{thm:known_landau} holds for $\sum n^\rho$, and $|n^{\rho}-\sqrt n\,n^{i\gamma}|\le n\log n\,|\beta-\frac12|$. Hypothesis~\ref{hyp:L} therefore follows from $\sum_{0<\gamma\le t}|\beta-\frac12|=o(t)$, whereas the density estimates of Selberg give $O(t)$ \cite[\S9.24]{Titchmarsh1986}. Theorems~\ref{thm:identity_general}--\ref{thm:main} do not use Hypothesis~\ref{hyp:L}.
\end{remark}

\begin{lemma}\label{lem:clock_equidistribution}
For fixed $\beta\ne0$, $\sum_{k\in\IT}e^{i\beta\tau_k}\ll_\beta\log T$. Consequently, for fixed integers $m,n\ge2$,
\begin{equation}
\frac1{\#\IT}\sum_{k\in\IT}\sin(\tau_k\log m)\sin(\tau_k\log n)\longrightarrow\frac12\,\mathbf 1_{m=n}.
\label{eq:clock_equidistribution}
\end{equation}
\end{lemma}

\begin{proof}
We may take $\beta>0$. Let $f(k)=\beta\tau_k/2\pi$ for real $k$. By \eqref{eq:clock_derivative}, $f'(k)=\beta/\log(\tau_k/2\pi)$ is positive and decreasing, and lies in $[\beta/\log(T/\pi),\frac12]$ for $k\in\IT$ and $T$ large; hence $\|f'\|\ge\beta/\log(T/\pi)$ and Theorem~\ref{thm:known_KL} gives the bound. For \eqref{eq:clock_equidistribution} write $2\sin a\sin b=\cos(a-b)-\cos(a+b)$ and use the bound with $\beta=\log(m/n)$ and $\beta=\log(mn)$, together with $\#\IT\asymp T\log T$.
\end{proof}

\begin{theorem}\label{thm:projection}
Assume Hypothesis~\ref{hyp:L}. For every fixed integer $n\ge2$,
\begin{equation}
\frac1{\#\IT}\sum_{k\in\IT}s_k\sin(\gamma_k\log n)\longrightarrow\frac{a_n}2=-\frac{\Lambda(n)}{2\pi\sqrt n\log n},
\label{eq:projection}
\end{equation}
and the same holds with $\tau_k$ in place of $\gamma_k$.
\end{theorem}

\begin{proof}
Let $\alpha=\log n$ and $\Phi(t)=\int_{14}^tS(u)n^{iu}du$. With $g(u)=n^{iu}/(i\alpha)$ and $S=N-\NO-\omega/\pi$, integration by parts gives
\begin{equation}
\Phi(t)=\bigl[Sg\bigr]_{14}^t-\frac1{i\alpha}\sum_{14<\gamma\le t}n^{i\gamma}+\int_{14}^tg\Bigl(\lL+\frac{\omega'}\pi\Bigr)du=\frac i\alpha\sum_{0<\gamma\le t}n^{i\gamma}+O(\log t),
\label{eq:projection_Phi}
\end{equation}
because $\int_{14}^tg\lL\,du=O(\log t)$ by a second integration by parts and $\omega'=O(u^{-2})$. Hypothesis~\ref{hyp:L} gives $\Phi(t)=-i\,\frac{\Lambda(n)}{2\pi\sqrt n\,\alpha}\,t+o(t)$. By partial summation, using $\frac{d}{dt}(t\lL(t))=\lL(t)+\frac1{2\pi}$,
\begin{equation}
\int_T^{2T}S(t)n^{it}\lL(t)\,dt=\bigl[\Phi\lL\bigr]_T^{2T}-\int_T^{2T}\frac{\Phi(t)}{2\pi t}\,dt=-i\,\frac{\Lambda(n)}{2\pi\sqrt n\,\alpha}\int_T^{2T}\lL\,dt+o(T\log T).
\label{eq:projection_partial}
\end{equation}
Taking imaginary parts, $\int_T^{2T}S(t)\sin(\alpha t)\lL(t)\,dt=\frac{a_n}2\int_T^{2T}\lL\,dt+o(T\log T)$. Now apply Lemma~\ref{lem:first_moment_identity} with $f(t)=\sin(\alpha t)$, $a=T$, $b=2T$ (we may assume that $T$ and $2T$ are not ordinates). The term with $\omega$ is $O(\log T)$, the boundary term is $O(\log^2T)$ and $\frac12\int_T^{2T}G^2|f'|\,dt\ll\int_T^{2T}S^2dt+O(1)\ll T\log\log T=o(T\log T)$ by Theorem~\ref{thm:known_selberg_moments}. Dividing by $\#\IT=\int_T^{2T}\lL\,dt+O(\log T)$ gives \eqref{eq:projection}. Finally $|\sin(\alpha\tau_k)-\sin(\alpha\gamma_k)|\le\alpha|\delta_k|$ and, by \eqref{eq:s_delta} and Theorem~\ref{thm:clock_moments}, $\sum_{k\in\IT}|s_k\delta_k|\ll(\log T)^{-1}\sum_{k\in\IT}s_k^2\ll\#\IT\,\log\log T/\log T$.
\end{proof}

\begin{corollary}\label{cor:gain}
Assume Hypothesis~\ref{hyp:L}. Let $A$ be a finite set of prime powers and $\lambda\in\R$. Then
\begin{equation}
\frac1{\#\IT}\sum_{k\in\IT}\Bigl(s_k^2-\bigl(s_k-\lambda S_A(\tau_k)\bigr)^2\Bigr)\longrightarrow(2\lambda-\lambda^2)\,v_A ,
\label{eq:gain}
\end{equation}
with $v_A$ as in \eqref{eq:SA_def}. The limit is maximal for $\lambda=1$, where it equals $v_A$.
\end{corollary}

\begin{proof}
The left side is $2\lambda\langle s_kS_A(\tau_k)\rangle-\lambda^2\langle S_A(\tau_k)^2\rangle$, where $\langle\cdot\rangle$ is the mean over $\IT$. By Theorem~\ref{thm:projection}, $\langle s_kS_A(\tau_k)\rangle\to\sum_{n\in A}a_n^2/2=v_A$, and by Lemma~\ref{lem:clock_equidistribution}, $\langle S_A(\tau_k)^2\rangle\to v_A$.
\end{proof}

\begin{remark}\label{rem:gain_full_factors}
For a finite set of primes $P$ and $A_{P,M}=\{p^m:p\in P,m\le M\}$ one has $v_{A_{P,M}}\uparrow c_P$ as $M\to\infty$ by \eqref{eq:cP_def}. The experiments of Section~\ref{sec:selected_computations} use the full local factors ($M=\infty$); the passage $M\to\infty$ inside the limit \eqref{eq:gain} is not covered by the proof, because the convergence in Theorem~\ref{thm:projection} is not shown to be uniform in $n$. Under the Riemann hypothesis the uniform version of Landau's formula \cite{Gonek1993} is the natural tool for that step; we do not carry it out here.
\end{remark}

\begin{remark}\label{rem:gain_meaning}
Corollary~\ref{cor:gain} shows that, among the multiples $\lambda S_A$ of the Euler phase, the Euler phase itself ($\lambda=1$) is optimal for the clock, and that its gain is the diagonal variance $v_A=\frac12\sum_{n\in A}a_n^2$. This is the full least-squares gain for those fixed components under the hypotheses used; the factor $1/2$ comes from the mean of $\sin^2$, not from an additional loss of half the predictable variance. Since $v_{A_X}\to\infty$ as $X\to\infty$ while $v_A$ is finite for every fixed $A$, a fixed set of primes recovers only a bounded part of the mean square $\frac1{2\pi^2}\log\log T$ of $s_k$, in agreement with Corollary~\ref{cor:fixed_cutoff}.
\end{remark}

\section{The one-step corrected clock}
\label{sec:one_step_clock}

The experiments evaluate the Euler phase at the clock point $\tau_k$, which is known, rather than at the unknown $\gamma_k$. The following identity relates the two.

\begin{proposition}\label{prop:one_step_decomposition}
Let $r_k(X)=\lL(\tau_k)\,\delta_k^X$ be the residual of the one-step clock in ordinal units. Then
\begin{equation}
r_k(X)=u_k(X)+\bigl(S_X(\gamma_k)-S_X(\tau_k)\bigr)+\eps_k,\qquad|\eps_k|\le\frac{\delta_k^2}{2\pi\min(\gamma_k,\tau_k)},
\label{eq:one_step_decomposition}
\end{equation}
and $|S_X(\gamma_k)-S_X(\tau_k)|\le M_X|\delta_k|$.
\end{proposition}

\begin{proof}
By \eqref{eq:one_step_clock}, $r_k(X)=\lL(\tau_k)\delta_k-S_X(\tau_k)$, and by \eqref{eq:s_delta}, $\lL(\tau_k)\delta_k=s_k+(\lL(\tau_k)-\lL(\xi_k))\delta_k$ with $|\lL(\tau_k)-\lL(\xi_k)|\le|\delta_k|/(2\pi\min(\gamma_k,\tau_k))$. Insert $s_k=u_k(X)+S_X(\gamma_k)$.
\end{proof}

\subsection{Two response functions and an exact transport identity}
\label{subsec:two_responses}
For a finite nonempty set of indices $I$, write $\langle h\rangle_I=|I|^{-1}\sum_{k\in I}h_k$. Unless a window is specified otherwise, $I=\IT$. For a prime power $n$, so that $a_n\ne0$, define
\begin{equation}
 \rho_n^\gamma(I)=\frac{2}{a_n}\langle s_k\sin(\gamma_k\log n)\rangle_I,
 \qquad
 \rho_n^\tau(I)=\frac{2}{a_n}\langle s_k\sin(\tau_k\log n)\rangle_I.
 \label{eq:two_responses}
\end{equation}
For a prime $p$, both normalizations are $-2\pi\sqrt p$ times the corresponding projection. These are finite-window statistics, defined without RH. A universal function of $\theta=\log p/\log(T/2\pi)$ would describe a limit of these statistics with $p$ varying with $T$; its existence is an additional assertion. Theorem~\ref{thm:projection} concerns fixed $n$, for which $\theta\to0$.

The distinction between the two evaluation points has the exact expression
\begin{equation}
 \begin{split}
 \rho_n^\gamma(I)-\rho_n^\tau(I)
 =-\frac{2\log n}{a_n}\int_0^1
 \left\langle s_k\delta_k
 \cos\bigl((\tau_k-v\delta_k)\log n\bigr)\right\rangle_I\,dv .
 \end{split}
 \label{eq:response_transport}
\end{equation}
This follows by differentiating the sine along the segment from $\tau_k$ to $\gamma_k=\tau_k-\delta_k$. The bounded-displacement interpolation used in the general RvM construction therefore also transports these observables, but its existence does not evaluate their transport integral. For fixed $n$, its absolute value tends to zero by Theorem~\ref{thm:clock_moments}, as in the proof of Theorem~\ref{thm:projection}. For $\log p\asymp\log T$, neither that argument nor a small mean square of $\delta_k$ makes the phase change $\delta_k\log p$ small. A law for $\rho^\gamma$ cannot consequently be substituted for a law for $\rho^\tau$.

\subsection{Exact accounting of the mean-square gain}
\label{subsec:exact_accounting}
Let $A$ be a finite set of prime powers and put
\begin{equation}
 C_A(I)=\langle S_A(\tau_k)^2\rangle_I-\frac12\sum_{n\in A}a_n^2.
 \label{eq:CA_def}
\end{equation}
It includes both the finite-window diagonal discrepancies and all cross terms. It need not have either sign.

\begin{proposition}[Finite-window gain]\label{prop:exact_gain}
For every such $I,A$ and every real $\lambda$,
\begin{equation}
 \begin{split}
 \left\langle s_k^2-(s_k-\lambda S_A(\tau_k))^2\right\rangle_I
 =\sum_{n\in A}a_n^2\left(\lambda\rho_n^\tau(I)-\frac{\lambda^2}{2}\right)
 -\lambda^2C_A(I).
 \end{split}
 \label{eq:exact_gain}
\end{equation}
The identity also holds for the full local factors $A=A_X$, with absolutely convergent sums; write $C_X=C_{A_X}$ in this case.
\end{proposition}
\begin{proof}
Expand the square and use $\langle s_kS_A(\tau_k)\rangle_I=\frac12\sum_{n\in A}a_n^2\rho_n^\tau(I)$. For fixed $X$, $\sum_{n\in A_X}|a_n|<\infty$, and $s_k$ is bounded on the finite set $I$, so the same argument passes to the full factors by absolute convergence.
\end{proof}
In particular the unit-coefficient gain is
\begin{equation}
 \begin{split}
 \langle s_k^2-(s_k-S_X(\tau_k))^2\rangle_I
 ={}&\sum_{p\le X}a_p^2\left(\rho_p^\tau-\tfrac12\right)-C_X\\
 &+\sum_{\substack{p\le X\\m\ge2}}a_{p^m}^2
       \left(\rho_{p^m}^\tau-\tfrac12\right).
 \end{split}
 \label{eq:full_factor_gain}
\end{equation}
For the prime-only phase the last sum is absent and $C_X$ must be computed for that phase. The two conventions must not be mixed. Formula~\eqref{eq:full_factor_gain} supplies an exact criterion for the sign of the gain without independence assumptions. A threshold $\rho_p^\tau=1/2$ describes a single diagonal prime contribution; it does not determine the optimum of a full Euler phase when $C_X$ and prime powers matter.

\subsection{The matrix projection and reconstruction of an ordinate}
\label{subsec:matrix_projection}
For a finite set $A$ of frequencies $n\ge2$, not necessarily prime powers, let
\begin{equation}
 \phi_n(k)=\sin(\tau_k\log n),\qquad
 G_{mn}=\langle\phi_m\phi_n\rangle_I,\qquad
 b_n=\langle s_k\phi_n(k)\rangle_I.
 \label{eq:Gram_projection}
\end{equation}
\begin{proposition}[Finite-window optimal coefficients]\label{prop:matrix_projection}
For $P_{\mathbf c}(k)=\sum_{n\in A}c_n\phi_n(k)$,
\begin{equation}
 \langle(s_k-P_{\mathbf c}(k))^2\rangle_I
 =\langle s_k^2\rangle_I-2\mathbf c^{\mathsf T}\mathbf b
       +\mathbf c^{\mathsf T}G\mathbf c.
 \label{eq:projection_quadratic}
\end{equation}
If $G$ is invertible, the unique minimizer is $\mathbf c=G^{-1}\mathbf b$ and the gain is $\mathbf b^{\mathsf T}G^{-1}\mathbf b$. In general the minimizers solve $G\mathbf c=\mathbf b$, with minimal-norm solution $G^\dagger\mathbf b$ and gain $\mathbf b^{\mathsf T}G^\dagger\mathbf b$.
\end{proposition}
\begin{proof}
Expansion gives \eqref{eq:projection_quadratic}. The matrix $G$ is positive semidefinite. If $G\mathbf v=0$, then $\langle P_{\mathbf v}^2\rangle_I=0$, so $P_{\mathbf v}=0$ on $I$ and $\mathbf v^{\mathsf T}\mathbf b=0$. Thus $\mathbf b$ belongs to the range of $G$. Completing the square on this range proves all assertions.
\end{proof}
For prime powers, $b_n=a_n\rho_n^\tau/2$. The approximation $G\approx\frac12I$ would give $c_n\approx a_n\rho_n^\tau$, but Lemma~\ref{lem:clock_equidistribution} establishes the limit only for fixed finite sets, not for sets growing with $T$.

Put $q_k=\lL(\tau_k)\delta_k=s_k+\eps_k$, with the error of Proposition~\ref{prop:one_step_decomposition}. Any chosen coefficient vector defines
\begin{equation}
 \widehat\gamma_k=\tau_k-\frac{P_{\mathbf c}(k)}{\lL(\tau_k)},\qquad
 \widehat\gamma_k-\gamma_k
   =\frac{s_k-P_{\mathbf c}(k)+\eps_k}{\lL(\tau_k)}.
 \label{eq:matrix_reconstruction}
\end{equation}
If diagonal prime coefficients $c_p=a_p\rho_p^\tau$ are used, the first formula reads
\begin{equation}
 \widehat\gamma_k=\tau_k+
 \frac{2}{\log(\tau_k/2\pi)}\sum_{p\le X}
 \rho_p^\tau\,\frac{\sin(\tau_k\log p)}{\sqrt p}.
 \label{eq:rho_tau_predictor}
\end{equation}
The response coefficients in this formula must be supplied by a proved law or by independent training data. Computing them from the same unknown ordinates produces a fitted reconstruction, not an independent prediction.

For exact least squares in ordinate units, the relevant matrix and vector instead are
\begin{equation}
 G^{(t)}_{mn}=\left\langle\frac{\phi_m\phi_n}{\lL(\tau_k)^2}\right\rangle_I,
 \qquad
 b^{(t)}_n=\left\langle\frac{\delta_k\phi_n}{\lL(\tau_k)}\right\rangle_I,
 \label{eq:weighted_Gram}
\end{equation}
and $G^{(t)}\mathbf c=\mathbf b^{(t)}$. This distinguishes the objective $\langle(\widehat\gamma-\gamma)^2\rangle_I$ from the ordinal objective in \eqref{eq:projection_quadratic}. Neither least-squares identity gives a pointwise bound uniform in $k$. Such a bound requires control of the individual residual in \eqref{eq:matrix_reconstruction}.

\subsection{Composite-frequency diagnostic}
\label{subsec:composite_diagnostic}
For every integer $n\ge2$ define the raw diagnostic normalization
\begin{equation}
\rho_n=-2\pi\sqrt n\langle s_k\sin(\tau_k\log n)\rangle.
\label{eq:raw_response}
\end{equation} It equals $\rho_p^\tau$ at primes, but for a prime power $n=p^m$ it equals $\rho_{p^m}^\tau/m$. The raw normalization is used in the composite-frequency formula and its numerical table; it is distinct from the unit-limit prime-power normalization in \eqref{eq:two_responses}.


\begin{heuristic}\label{heur:second_order}
This is a heuristic argument. In Theorem~\ref{thm:projection} the term $-\frac12\int G^2f'$ of \eqref{eq:first_moment_identity} and the replacement of $\gamma_k$ by $\tau_k$ are of relative size $1/\log T$. Modelling $S$ by $\sum a_m\sin(t\log m)$, the coefficient of $\cos(t\log n)$ in $S^2$ for $n=p^aq^b$ with distinct primes $p,q$ is $-a_{p^a}a_{q^b}$, and the two corrections combine to the prediction
\begin{equation}
\rho_n:=-2\pi\sqrt n\,\Bigl\langle s_k\sin(\tau_k\log n)\Bigr\rangle\approx\frac{\Lambda(n)}{\log n}+\frac{\log n}{ab\,\log(T/2\pi)}\,\mathbf 1_{n=p^aq^b},
\label{eq:second_order}
\end{equation}
the first term being the limit given by Theorem~\ref{thm:projection}.
\end{heuristic}

\section{A moving-prime response law and its exact arithmetic target}
\label{sec:moving_response}
The function singled out by the prime-resonant calculation is
\begin{equation}
 \Phi(\theta)=(1+\theta)\log(1+\theta)+(1-\theta)\log(1-\theta),
 \qquad R(\theta)=1-\Phi(\theta),\quad 0<\theta<1.
 \label{eq:universal_candidate}
\end{equation}
The calculation in Section~\ref{sec:sine_response} proves its appearance in a specified determinantal response model. Its identification with the responses of the zeta ordinates is the arithmetic problem formulated below. No theorem for the zeta ordinates is inferred merely from a GUE interpretation.

\subsection{An exact formula for the response at a genuine ordinate}
\label{subsec:gamma_response_identity}
Let $T,2T$ be points which are not zero ordinates, put $M_T=\#\IT$, $L=\log(T/2\pi)$, $G=N-\NO=S+\omega/\pi$, and, for a prime $p$, put $w=\log p$ and $\theta=w/L$. Define
\begin{equation}
 A_{T,p}=-\frac{2\pi\sqrt p}{M_T}\int_T^{2T}G(t)\sin(wt)\lL(t)\,dt,
 \qquad C_{T,p}=\frac1T\int_T^{2T}S(t)^2p^{it}\,dt.
 \label{eq:ATp_CTp}
\end{equation}
\begin{proposition}[Separation of one- and two-point terms]\label{prop:gamma_response_exact}
For every such $T,p$,
\begin{equation}
 \begin{split}
 \rho_p^\gamma(\IT)={}& A_{T,p}
 +\frac{\pi\sqrt p\,wT}{M_T}\Re C_{T,p}
 -\frac{\pi\sqrt p}{M_T}[G(t)^2\sin(wt)]_T^{2T}\\
 &+\frac{\pi\sqrt p\,w}{M_T}
       \int_T^{2T}(G(t)^2-S(t)^2)\cos(wt)\,dt.
 \end{split}
 \label{eq:gamma_response_exact}
\end{equation}
In particular, if $p=(T/2\pi)^{\theta+o(1)}$ and $\theta$ stays in a compact subset of $(0,1)$, then
\begin{equation}
 \rho_p^\gamma=A_{T,p}
 +2\pi^2\theta\sqrt p\,\bigl(1+O(1/L)\bigr)\Re C_{T,p}+o(1).
 \label{eq:gamma_response_scaled}
\end{equation}
The factor $1+O(1/L)$ is retained until the size of $\sqrt p\,C_{T,p}$ is controlled.
\end{proposition}
\begin{proof}
Apply Lemma~\ref{lem:first_moment_identity} with $f(t)=\sin(wt)$, and multiply by $-2\pi\sqrt p/M_T$. This gives \eqref{eq:gamma_response_exact}, including its boundary sign. Since $\omega(t)=O(1/T)$ and $S(t)=O(\log T)$ on the window, $\int_T^{2T}|G^2-S^2|\,dt\ll\log T$. The boundary contribution is $O(\sqrt p\log T/T)$; the last integral contributes $O(\sqrt p\,w/T)$. Both tend to zero in the stated range. Finally $M_T=TL/(2\pi)+O(T)$, which gives \eqref{eq:gamma_response_scaled}.
\end{proof}

\begin{hypothesis}[Uniform one-particle input]\label{hyp:uniform_prime}
For every fixed $0<\theta_0<\theta_1<1$, uniformly over primes satisfying $\theta_0\le\log p/L\le\theta_1$, one has $A_{T,p}=1+o(1)$ as $T\to\infty$.
\end{hypothesis}
This is a moving-prime input, distinct from Hypothesis~\ref{hyp:L}. One sufficient, stronger explicit-formula input is
\begin{equation}
 \sup_{T\le t\le2T}
 \left|\sum_{T<\gamma_k\le t}p^{i\gamma_k}
       +\frac{(t-T)\log p}{2\pi\sqrt p}\right|
   =o\left(\frac{T\log p}{\sqrt p}\right)
 \label{eq:uniform_Landau_sufficient}
\end{equation}
uniformly in that range. Indeed, integrate $G(t)e^{iwt}\lL(t)$ by parts with the primitive $e^{iwt}\lL(t)/(iw)$; the extra term containing $\lL'(t)$ is $O(\log T/w)$ and the endpoint terms are $O(\log^2T/w)$. A second partial summation applied to \eqref{eq:uniform_Landau_sufficient} gives the main term $-i(2\pi\sqrt p)^{-1}\int_T^{2T}\lL(t)\,dt$ and an error $o(M_T/\sqrt p)$. This proves the sufficient implication. The fixed-base formula of Theorem~\ref{thm:known_landau} cannot be used to assert this uniform estimate. Unconditionally it concerns $p^\rho$, not $\sqrt p\,p^{i\gamma}$; that replacement requires RH or a separate quantitative estimate for the real parts of the zeros.

\subsection{The scalar arithmetic assertion}
\label{subsec:scalar_response_target}
The prime-response target is
\begin{equation}
 C_{T,p}=-\frac{1}{2\pi^2\sqrt p}\frac{\Phi(\theta)}{\theta}
                    +o(p^{-1/2}).
 \label{eq:CTp_target}
\end{equation}
It is a complex assertion: it specifies the real main term and requires $\Im C_{T,p}=o(p^{-1/2})$. For the real statistic $\rho_p^\gamma$, the corresponding real-part assertion suffices.

\begin{proposition}[Conditional moving-prime law]\label{prop:conditional_universal}
Assume Hypothesis~\ref{hyp:uniform_prime} and the real part of \eqref{eq:CTp_target}, uniformly for $\theta$ in a compact subset of $(0,1)$. Then
\begin{equation}
 \rho_p^\gamma(\IT)=1-\Phi(\theta)+o(1).
 \label{eq:conditional_universal}
\end{equation}
Conversely, under the same one-particle input, \eqref{eq:conditional_universal} is equivalent to that real-part assertion. Neither implication asserts a formula for $\rho_p^\tau$.
\end{proposition}
\begin{proof}
Use \eqref{eq:gamma_response_scaled}; the real-part assertion bounds $\sqrt p\,\Re C_{T,p}$, so the relative $O(1/L)$ factor contributes $o(1)$. Conversely solve that exact coefficient relation for $\Re C_{T,p}$.
\end{proof}
For $|\theta|<1$,
\begin{equation}
 \Phi(\theta)=\theta^2+\frac{\theta^4}{6}+\frac{\theta^6}{15}
                    +\frac{\theta^8}{28}+\cdots
 =2\sum_{j\ge1}\frac{\theta^{2j}}{(2j-1)(2j)}.
 \label{eq:Phi_series}
\end{equation}
Thus $1-\theta^2$ is only the quadratic truncation of the candidate law. All higher terms require arithmetic information at the same scale as \eqref{eq:CTp_target}.

\section{A Fourier reduction of the prime target}
\label{sec:Fourier_reduction}
\subsection{The distributional measure and the shifted correlation}
\label{subsec:distributional_measure}
Put $H(t)=\NO(t)+\omega(t)/\pi$ and use the convention $\widehat f(\xi)=\int_\R f(t)e^{i\xi t}\,dt$. Define
\begin{equation}
 f_T=S\mathbf1_{[T,2T]},\qquad \mu_T=Df_T,
 \qquad Z_T(\xi)=\int_\R e^{i\xi t}\,d\mu_T(t).
 \label{eq:measure_def}
\end{equation}
\begin{theorem}[Shifted two-point representation]\label{thm:Fourier_exact}
Without RH or a simplicity assumption,
\begin{equation}
 \begin{split}
 d\mu_T={}&\mathbf1_{(T,2T)}(dN-H'(t)\,dt)
          +S(T)\delta_T-S(2T)\delta_{2T},\\
 Z_T(\xi)={}&\sum_{T<\gamma<2T}\nu(\gamma)e^{i\xi\gamma}
       -\int_T^{2T}H'(t)e^{i\xi t}\,dt
       +S(T)e^{i\xi T}-S(2T)e^{2i\xi T},
 \end{split}
 \label{eq:measure_expansion}
\end{equation}
where $\nu(\gamma)$ is the number of zeros with that ordinate, counted with their analytic multiplicities. Moreover,
\begin{equation}
 Z_T(0)=0,\qquad Z_T(\xi)=-i\xi\widehat f_T(\xi),
 \qquad \frac{Z_T(\xi)}{\xi}=-i\widehat f_T(\xi),
 \label{eq:ZT_removable}
\end{equation}
the last quotient being continued at zero. For every real $w$,
\begin{equation}
 \int_T^{2T}S(t)^2e^{iwt}\,dt
  =-\frac1{2\pi}\int_\R
       \frac{Z_T(\xi)Z_T(w-\xi)}{\xi(w-\xi)}\,d\xi .
 \label{eq:Fourier_convolution}
\end{equation}
The integral, with the continued quotients, is absolutely convergent.
\end{theorem}
\begin{proof}
Between ordinates $DS=dN-H'\,dt$. Differentiating the cutoff supplies exactly the two endpoint atoms of \eqref{eq:measure_expansion}. Integration by parts gives $\widehat{Df_T}=-i\xi\widehat f_T$, including $Z_T(0)=0$. The function $f_T$ is compactly supported and belongs to $L^1\cap L^2$. The product-convolution identity therefore yields \eqref{eq:Fourier_convolution}. Absolute convergence follows from Cauchy--Schwarz and Plancherel, since
\begin{equation}
 \int_\R|\widehat f_T(\xi)\widehat f_T(w-\xi)|\,d\xi
                   \le\|\widehat f_T\|_2^2<\infty.
 \label{eq:Fourier_absolute}
\end{equation}
All identities use the full measure $dN$, so repeated ordinates cause no difficulty.
\end{proof}
With $w=\theta L=\log p$ and $\xi=\alpha L$, define
\begin{equation}
 \mathcal B_T(\theta)=\int_\R
 \frac{Z_T(\alpha L)Z_T((\theta-\alpha)L)}{\alpha(\theta-\alpha)}\,d\alpha.
 \qquad C_{T,p}=-\frac{\mathcal B_T(\theta)}{2\pi TL}.
 \label{eq:BT_scaled}
\end{equation}
Consequently the complex assertion \eqref{eq:CTp_target} is exactly equivalent to
\begin{equation}
 \mathcal B_T(\theta)=\frac{TL}{\pi\sqrt p}\frac{\Phi(\theta)}{\theta}
                              +o(TL/\sqrt p).
 \label{eq:BT_target}
\end{equation}
Taking only real parts gives an equivalence only for real parts. The additional assertion $\Im\mathcal B_T=o(TL/\sqrt p)$ is necessary for the complex target.

Expanding the two zero sums produces factors
\begin{equation}
 e^{i\alpha L\gamma}e^{i(\theta-\alpha)L\gamma'}
        =p^{i\gamma'}e^{i\alpha L(\gamma-\gamma')}.
 \label{eq:shifted_pair_weight}
\end{equation}
The observable is therefore a two-point correlation shifted by a nonzero prime frequency. An ordinary stationary pair correlation, depending only on $\gamma-\gamma'$, does not determine it. At $\theta=0$ the product becomes $|Z_T(\alpha L)|^2$; that special case does not supply the information at $\theta>0$.

\subsection{Smooth windows and the arithmetic scale}
\label{subsec:smooth_windows}
For real $w_0\in C_c^\infty((1,2))$, put $f_{T,w_0}(t)=w_0(t/T)S(t)$. The exact replacement for the first line of \eqref{eq:measure_expansion} is
\begin{equation}
 Df_{T,w_0}=w_0(t/T)(dN-H'(t)\,dt)
             +T^{-1}w_0'(t/T)S(t)\,dt.
 \label{eq:smooth_measure}
\end{equation}
The last term is part of the measure and cannot be omitted. Applying Theorem~\ref{thm:Fourier_exact} to this function represents $\int w_0(t/T)^2S(t)^2e^{iwt}\,dt$. If the diagonal normalization is to be unchanged, use $\int_\R w_0(u)^2\,du=1$. A sharp-window conjecture and a smooth-window conjecture require their own quantitative transfer; they are not silently identified here.

There is an elementary reason that the prime polynomial alone does not settle \eqref{eq:CTp_target}. For the finite polynomial
\begin{equation}
 P_Y(t)=-\frac1\pi\sum_{p\le Y}\frac{\sin(t\log p)}{\sqrt p},
 \label{eq:PY_audit}
\end{equation}
the square has no exact Fourier frequency $\log p$ for a prime $p$. If prime powers are included, the resonances $p^a,p^{a+1}$ first have size $p^{-3/2}$. In contrast, the $q=p$ term of $P_Y'$ contributes
\begin{equation}
 \frac1T\int_\R W(t/T)P_Y'(t)p^{it}\,dt
       =-\frac{\log p}{2\pi\sqrt p}+\mathcal E_{T,p,Y},
 \label{eq:derivative_diagonal}
\end{equation}
when $p\le Y$ and $\int W=1$. Here $\mathcal E_{T,p,Y}$ is the explicit sum of all nonzero-frequency Fourier transforms of $W$; its smallness for moving $p,Y$ needs a uniform separation estimate. Equations~\eqref{eq:PY_audit}--\eqref{eq:derivative_diagonal} identify exact resonances, not an unproved uniform estimate for the remaining terms. After division by $\pi L$, the displayed diagonal is $-\theta/(2\pi^2\sqrt p)$, the linear term of $-\Phi(\theta)/(2\pi^2\sqrt p\,\theta)$.

Writing $S=P_Y+R_Y$ gives $S^2=P_Y^2+2P_YR_Y+R_Y^2$. Selberg's estimate for the appropriate short polynomial gives $T^{-1}\int|R_Y|^2\ll1$, and hence only
\begin{equation}
 \left|T^{-1}\int_T^{2T}R_Y(t)^2p^{it}\,dt\right|\ll1.
 \label{eq:Selberg_scale_gap}
\end{equation}
The desired error is $o(p^{-1/2})=T^{-\theta/2+o(1)}$. The corresponding mixed term also needs oscillatory cancellation. Additional absolute moments by themselves do not bridge this difference in scale. Theorems~\ref{thm:identity_general} and~\ref{thm:main} remain unconditional throughout this reduction.

\section{A determinantal model for the response function}
\label{sec:sine_response}
\subsection{A finite model and a real exponential perturbation}
\label{subsec:tilt_definition}
To define the perturbation without infinite-volume or complex-probability ambiguities, start on a circle of length $R_0$. Let $P_N$ be the orthogonal projection onto $N$ consecutive Fourier modes and use its determinantal point process, which has exactly $N$ points. For $q=j/R_0$, let $X_q$ be the centred linear statistic of $e^{2\pi iqx}$. In particular $X_0=0$ and $\overline{X_q}=X_{-q}$. Choose $N/R_0\to1$; the local projection kernel tends to the unit-density sine kernel $\sin(\pi(x-y))/(\pi(x-y))$.

For $\theta>0$ on the frequency lattice define an actual probability measure by
\begin{equation}
 d\mathbb P_{u,v}=\frac{\exp\{-u\Re X_{-\theta}-v\Im X_{-\theta}\}}
 {\mathbb E\exp\{-u\Re X_{-\theta}-v\Im X_{-\theta}\}}\,d\mathbb P_0,
 \qquad u,v\in\R.
 \label{eq:real_tilt}
\end{equation}
The statistics are bounded in this finite model, so differentiation at $(0,0)$ is justified. For the covariance without complex conjugation,
\begin{equation}
 (\partial_u+i\partial_v)
 \operatorname{Cov}_{u,v}(X_\alpha,X_{\theta-\alpha})\big|_{(0,0)}
   =-\kappa_3(X_\alpha,X_{\theta-\alpha},X_{-\theta}).
 \label{eq:tilt_cumulant}
\end{equation}
This follows by differentiating the normalized expectation and subtracting the derivatives of its means. Thus a specified exponential perturbation produces the third connected cumulant. Matching one Fourier coefficient of a point process would not, by itself, determine this perturbation.

\subsection{Evaluation of the connected kernel}
\label{subsec:connected_kernel}
Write $x_+=\max(x,0)$ and $h(q)=(1-|q|)_+$.
\begin{theorem}[Connected response kernel]\label{thm:sine_kernel}
Let the lattice frequencies tend to $\alpha,\theta-\alpha,-\theta$, with sum zero, where $0<\theta<1$. In the limit $R_0,N\to\infty$, $N/R_0\to1$,
\begin{equation}
 \frac1{R_0}\kappa_3(X_\alpha,X_{\theta-\alpha},X_{-\theta})
 \longrightarrow k_\theta(\alpha)
 =1-h(\alpha)-h(\theta-\alpha)-h(\theta)
       +2(1-M)_+,
 \label{eq:kernel_intersections}
\end{equation}
where $M=\max(|\alpha|,|\theta-\alpha|,\theta)$. Equivalently,
\begin{equation}
 k_\theta(\alpha)=
 \begin{cases}
 \theta,&\alpha\le-1,\\
 \theta-1-\alpha,&-1\le\alpha\le\theta-1,\\
 0,&\theta-1\le\alpha\le1,\\
 \alpha-1,&1\le\alpha\le1+\theta,\\
 \theta,&\alpha\ge1+\theta.
 \end{cases}
 \label{eq:kernel_piecewise}
\end{equation}
\end{theorem}
\begin{proof}
For commuting multiplication operators $f,g,h$, expansion of the determinantal generating function $\det(I+P_N(e^{zf+ug+vh}-1)P_N)$ gives
\begin{equation}
 \begin{split}
 \kappa_3(f,g,h)={}&\operatorname{Tr}(fghP_N)
 -\operatorname{Tr}(fgP_NhP_N)-\operatorname{Tr}(fhP_NgP_N)\\
 &-\operatorname{Tr}(ghP_NfP_N)
 +\operatorname{Tr}(fP_NgP_NhP_N)
 +\operatorname{Tr}(fP_NhP_NgP_N).
 \end{split}
 \label{eq:third_cumulant_trace}
\end{equation}
This is determinantal cumulant algebra, rather than a claim of a new general cumulant formula. Multiplication by $e^{2\pi ijx/R_0}$ shifts a Fourier mode by $j$. A trace of two projections therefore counts the overlap of two integer intervals of length $N$, and a trace of three counts their triple overlap. For three shifts with sum zero, the triple overlap has length $(N-\operatorname{range}(0,j_1,j_1+j_2))_+$. The two cyclic orders have the same range, whose rescaled limit is $M$. Divide \eqref{eq:third_cumulant_trace} by $R_0$ and use the interval-overlap limits. This gives \eqref{eq:kernel_intersections}; evaluating the positive parts on the five stated intervals gives \eqref{eq:kernel_piecewise}. If any frequency is zero, both the centred cumulant and the resulting kernel vanish, as required.
\end{proof}

\begin{proposition}[Integral and support]\label{prop:kernel_integral}
For $0<\theta<1$,
\begin{equation}
 \int_\R\frac{k_\theta(\alpha)}{\alpha(\theta-\alpha)}\,d\alpha
       =-\frac{2\Phi(\theta)}{\theta}.
 \label{eq:kernel_integral}
\end{equation}
The integral is ordinary and absolutely convergent. Moreover,
\begin{equation}
 |\alpha|+|\theta-\alpha|+\theta<2
       \quad\Longrightarrow\quad k_\theta(\alpha)=0.
 \label{eq:kernel_support_barrier}
\end{equation}
\end{proposition}
\begin{proof}
The kernel vanishes on a neighborhood of $0$ and $\theta$ and is bounded on the tails. Its symmetry under $\alpha\mapsto\theta-\alpha$ reduces the integral to twice its part on $[1,\infty)$. The tail integral on $[1+\theta,\infty)$ equals $-\log(1+\theta)$. The integral on $[1,1+\theta]$ equals
\begin{equation}
 -\frac{\log(1+\theta)}{\theta}
       +\frac{\theta-1}{\theta}\log(1-\theta).
 \label{eq:kernel_ramp_integral}
\end{equation}
Adding and doubling proves \eqref{eq:kernel_integral}. The nonzero kernel lies in $\alpha<\theta-1$ or $\alpha>1$. In both intervals the sum of absolute frequencies is greater than $2$, which proves \eqref{eq:kernel_support_barrier}, including vanishing at the boundary by continuity.
\end{proof}
Thus a test confined to the displayed restricted-support region sees zero of this connected response. This statement is proved for the model with the frequency normalization used here. It does not identify all published restricted-support theorems with one another and does not apply one outside its stated range.

\subsection{Normalization of the model and the missing arithmetic bridge}
\label{subsec:model_bridge}
The same projection calculation gives
\begin{equation}
 R_0^{-1}\kappa_2(X_\theta,X_{-\theta})\longrightarrow\theta,
 \qquad 0<\theta<1.
 \label{eq:sine_susceptibility}
\end{equation}
For a zeta window the unfolded length is $\mathscr L_T=TL/(2\pi)$. The one-particle prime amplitude in \eqref{eq:uniform_Landau_sufficient} is $-\mathscr L_T\theta/\sqrt p$. This matches the linear response of \eqref{eq:real_tilt} with intensity $p^{-1/2}$. It is a normalization of the model, not a characterization of the actual zeta measure. The corresponding proposed connected product is $-\mathscr L_Tk_\theta(\alpha)/\sqrt p$.

If this last expression could be substituted into the full weighted integral \eqref{eq:BT_scaled} with an error $o(TL/\sqrt p)$, Proposition~\ref{prop:kernel_integral} would give exactly \eqref{eq:BT_target} and hence the complex target \eqref{eq:CTp_target}. Compactly supported distributional convergence alone would not justify this substitution: the two removable singularities and the tails of the weighted integral also need uniform control.

\begin{hypothesis}[Weighted arithmetic bridge]\label{hyp:weighted_bridge}
With the sharp-window measure of \eqref{eq:measure_def}, uniformly for primes with $\theta$ in any compact subset of $(0,1)$,
\begin{equation}
 \int_\R\frac{Z_T(\alpha L)Z_T((\theta-\alpha)L)
             +\mathscr L_Tp^{-1/2}k_\theta(\alpha)}
                 {\alpha(\theta-\alpha)}\,d\alpha
                         =o(TL/\sqrt p).
 \label{eq:weighted_arithmetic_bridge}
\end{equation}
All quotients from $Z_T$ use their continued values. The kernel term vanishes near the singularities.
\end{hypothesis}
This hypothesis is precisely the complex scalar target after the proved integral evaluation; it is not asserted as a consequence of GUE or of RH. Under Hypotheses~\ref{hyp:uniform_prime} and~\ref{hyp:weighted_bridge}, Proposition~\ref{prop:conditional_universal} gives the law $R(\theta)$ for $\rho_p^\gamma$. The response kernel and its integral are proved model results; the arithmetic bridge, and a separate same-clock law for $\rho_p^\tau$, remain open.


\section{Selected computations}\label{sec:selected_computations}
The computations in this section are finite statements about the defined observables.  They are not used to infer an asymptotic cutoff exponent.\subsection{Direct defects and moving-prime responses}
\label{subsec:direct_extensions}
\begin{experiment}[Direct measurement of the theorem's observable]\label{exp:direct_defects}
For $T=2^r$, $r\in\{14,16,18,19\}$, and $c\in\{1/10,1/5,1/3,1/2\}$, we compute the full local-factor phase at the true ordinates, with one cutoff $X=T^c$ throughout each window, and form $u_k(X)=k-1/2-\NO(\gamma_k)-S_X(\gamma_k)$. The numbers of ordinates are respectively $21\,519$, $100\,536$, $459\,982$ and $977\,804$. The last window, entirely contained in the data, gives Table~\ref{tab:direct_top}; all sixteen measurements and the one-step comparison appear in Appendix~\ref{sec:selected_computations}. Exponents smaller than $1/10$ on the stated trial grid give $X<2$ at these heights and are excluded, consistently with the theorem's range. The finite means do not identify a limiting constant at fixed $c$.
\end{experiment}
\begin{table}[htbp]
\centering\small
\begin{tabular}{rrrrr}
\toprule
$c$ & $\#\{p\le X\}$ & $\langle u_k(X)^2\rangle$ & $\frac1{2\pi^2}\log(1/c)$ & $\langle r_k(X)^2\rangle$\\
\midrule
$1/10$ & 2 & 0.065599 & 0.116650 & 0.065517\\
$1/5$ & 6 & 0.040530 & 0.081535 & 0.040623\\
$1/3$ & 22 & 0.022911 & 0.055656 & 0.023446\\
$1/2$ & 128 & 0.012187 & 0.035115 & 0.014567\\
\bottomrule
\end{tabular}
\caption{Direct ordinal defects on $2^{19}<\gamma_k\le2^{20}$, with $977\,804$ ordinates and $X=(2^{19})^c$. The one-step residual has no saturation. These are different observables.}
\label{tab:direct_top}
\end{table}

\begin{experiment}[Individual-prime and grouped responses]\label{exp:moving_responses}
For every one of the $2262$ primes $p\le20000$ we compute both responses in \eqref{eq:two_responses} on the five ordinal windows after $k=1000$. We use the exact $s_k=k-1/2-\NO(\gamma_k)$, and assign to a prime the horizontal coordinate $\theta=\log p/\operatorname{median}_{k\in I}\log(\tau_k/2\pi)$. These are finite ordinal windows, rather than the dyadic height windows of the arithmetic hypotheses. Only primes with $0<\theta<1$ are compared with $R(\theta)$. The complete individual responses are retained in the companion CSV, and Table~\ref{tab:response_errors} distinguishes individual errors from averages in bins of width $0.05$, using only bins with at least five primes. The grouped comparison supports the shape of the candidate; it supplies neither a uniform individual-prime error of $0.03$ nor an asymptotic proof.
\end{experiment}
\begin{table}[htbp]
\centering\small
\begin{tabular}{lrrrr}
\toprule
ordinal window & primes with $\theta<1$ & max individual error & RMS error & max bin-mean error\\
\midrule
$(10^3,10^4]$ & 159 & 0.34557 & 0.09425 & 0.02637\\
$(10^4,10^5]$ & 901 & 0.44085 & 0.08511 & 0.02371\\
$(10^5,5\cdot10^5]$ & 2262 & 0.27857 & 0.06429 & 0.02972\\
$(5\cdot10^5,10^6]$ & 2262 & 0.31682 & 0.06086 & 0.03531\\
$(10^6,K]$ & 2262 & 0.16425 & 0.03309 & 0.01951\\
\bottomrule
\end{tabular}
\caption{Errors of $\rho_p^\gamma-R(\theta)$ in Experiment~\ref{exp:moving_responses}. Bin errors are absolute means of the signed individual errors, not maxima over individual primes.}
\label{tab:response_errors}
\end{table}
In the top window, the means of $\rho_p^\tau$ in the bins $[0.65,0.70)$, $[0.70,0.75)$ and $[0.75,0.80)$ are $0.66402$, $0.65240$ and $0.66417$. This finite-height plateau is different from $R(\theta)$ and does not establish a universal same-clock law. In particular, the prime near $p=17333$ has an individual $\rho_p^\gamma$ error about $0.16425$, despite the much smaller bin-mean discrepancies.


\section{Conclusion}
The identities, asymptotic estimates, conditional consequences and model calculations above have distinct hypotheses and are not merged into a single claim.  In particular, a bound for $u_k(X)$ is not a root-selection theorem unless the strict derivative hypothesis in Proposition~\ref{prop:localization} holds, and the determinantal calculation is not an arithmetic theorem about the zeta zeros.
\section*{Statements and Declarations}
\paragraph{Competing interests.} The author declares no competing interests.
\paragraph{Funding.} No external funding is declared.
\paragraph{Data and code availability.} The accompanying archive contains the Python programs, data-independent checks and output tables. The two large binary ordinate files are not redistributed.
\begin{thebibliography}{99}

\bibitem{BellottiWong2025}
C.~Bellotti and P.-J. Wong, Improved estimates for the argument and zero-counting function of the Riemann zeta-function, \emph{Math. Comp.} (2025, published online), \href{https://doi.org/10.1090/mcom/4133}{doi:10.1090/mcom/4133}; arXiv:2412.15470.

\bibitem{Billingsley1995}
P.~Billingsley, \emph{Probability and Measure}, 3rd ed., Wiley, New York, 1995. (No DOI.)

\bibitem{Corless1996}
R.~M. Corless, G.~H. Gonnet, D.~E.~G. Hare, D.~J. Jeffrey and D.~E. Knuth, On the Lambert $W$ function, \emph{Adv. Comput. Math.} \textbf{5} (1996), 329--359, \href{https://doi.org/10.1007/BF02124750}{doi:10.1007/BF02124750}.

\bibitem{FrancaLeClair2014}
G.~Fran\c{c}a and A.~LeClair, Statistical and other properties of Riemann zeros based on an explicit equation for the $n$-th zero on the critical line, preprint, arXiv:1307.8395 (2013), \href{https://doi.org/10.48550/arXiv.1307.8395}{doi:10.48550/arXiv.1307.8395}.

\bibitem{FrancaLeClair2015}
G.~Fran\c{c}a and A.~LeClair, Transcendental equations satisfied by the individual zeros of Riemann $\zeta$, Dirichlet and modular $L$-functions, \emph{Commun. Number Theory Phys.} \textbf{9} (2015), no.~1, 1--50, \href{https://doi.org/10.4310/CNTP.2015.v9.n1.a1}{doi:10.4310/CNTP.2015.v9.n1.a1}; arXiv:1502.06003.

\bibitem{Fujii1974}
A.~Fujii, On the zeros of Dirichlet $L$-functions. I, \emph{Trans. Amer. Math. Soc.} \textbf{196} (1974), 225--235, \href{https://doi.org/10.1090/S0002-9947-1974-0349603-2}{doi:10.1090/S0002-9947-1974-0349603-2}.

\bibitem{Gonek1993}
S.~M. Gonek, An explicit formula of Landau and its applications to the theory of the zeta-function, in: \emph{A Tribute to Emil Grosswald: Number Theory and Related Analysis}, Contemp. Math. \textbf{143}, Amer. Math. Soc., Providence, RI, 1993, 395--413, \href{https://doi.org/10.1090/conm/143/1210528}{doi:10.1090/conm/143/1210528}.

\bibitem{Gonek2012}
S.~M. Gonek, Finite Euler products and the Riemann hypothesis, \emph{Trans. Amer. Math. Soc.} \textbf{364} (2012), 2157--2191, \href{https://doi.org/10.1090/S0002-9947-2011-05546-7}{doi:10.1090/S0002-9947-2011-05546-7}; arXiv:0704.3448.

\bibitem{GrahamKolesnik1991}
S.~W. Graham and G.~Kolesnik, \emph{Van der Corput's Method of Exponential Sums}, London Math. Soc. Lecture Note Ser. \textbf{126}, Cambridge University Press, Cambridge, 1991, \href{https://doi.org/10.1017/CBO9780511661976}{doi:10.1017/CBO9780511661976}.

\bibitem{HasanalizadeShenWong2022}
E.~Hasanalizade, Q.~Shen and P.-J. Wong, Counting zeros of the Riemann zeta function, \emph{J. Number Theory} \textbf{235} (2022), 219--241, \href{https://doi.org/10.1016/j.jnt.2021.06.032}{doi:10.1016/j.jnt.2021.06.032}; arXiv:2107.06506.

\bibitem{Korolev2012}
M.~A. Korolev, On Gram's law in the theory of the Riemann zeta function, \emph{Izv. Math.} \textbf{76} (2012), no.~2, 275--309, \href{https://doi.org/10.1070/IM2012v076n02ABEH002583}{doi:10.1070/IM2012v076n02ABEH002583}; Russian original \emph{Izv. Ross. Akad. Nauk Ser. Mat.} \textbf{76} (2012), no.~2, 67--102, \href{https://doi.org/10.4213/im5558}{doi:10.4213/im5558}; arXiv:1011.3997.

\bibitem{Landau1912}
E.~Landau, \"Uber die Nullstellen der Zetafunktion, \emph{Math. Ann.} \textbf{71} (1912), 548--564, \href{https://doi.org/10.1007/BF01456808}{doi:10.1007/BF01456808}.

\bibitem{LeClair2021}
A.~LeClair, Riemann hypothesis and random walks: the zeta case, \emph{Symmetry} \textbf{13} (2021), no.~11, 2014, \href{https://doi.org/10.3390/sym13112014}{doi:10.3390/sym13112014}; arXiv:1601.00914.

\bibitem{Littlewood1924}
J.~E. Littlewood, On the zeros of the Riemann zeta-function, \emph{Proc. Cambridge Philos. Soc.} \textbf{22} (1924), 295--318, \href{https://doi.org/10.1017/S0305004100014225}{doi:10.1017/S0305004100014225}.

\bibitem{Mertens1874}
F.~Mertens, Ein Beitrag zur analytischen Zahlentheorie, \emph{J. Reine Angew. Math.} \textbf{78} (1874), 46--62, \href{https://doi.org/10.1515/crll.1874.78.46}{doi:10.1515/crll.1874.78.46}.

\bibitem{MontgomeryVaughan1974}
H.~L. Montgomery and R.~C. Vaughan, Hilbert's inequality, \emph{J. London Math. Soc. (2)} \textbf{8} (1974), 73--82, \href{https://doi.org/10.1112/jlms/s2-8.1.73}{doi:10.1112/jlms/s2-8.1.73}.

\bibitem{Selberg1946a}
A.~Selberg, Contributions to the theory of the Riemann zeta-function, \emph{Arch. Math. Naturvid.} \textbf{48} (1946), no.~5, 89--155. (No DOI.)

\bibitem{Selberg1946b}
A.~Selberg, Contributions to the theory of Dirichlet's $L$-functions, \emph{Skr. Norske Vid.-Akad. Oslo I} \textbf{1946}, no.~3, 2--62. (No DOI.)

\bibitem{Titchmarsh1986}
E.~C. Titchmarsh, \emph{The Theory of the Riemann Zeta-Function}, 2nd ed., revised by D.~R. Heath-Brown, Clarendon Press, Oxford, 1986. (No DOI.)

\bibitem{Trudgian2014}
T.~S. Trudgian, An improved upper bound for the argument of the Riemann zeta-function on the critical line II, \emph{J. Number Theory} \textbf{134} (2014), 280--292, \href{https://doi.org/10.1016/j.jnt.2013.07.017}{doi:10.1016/j.jnt.2013.07.017}; arXiv:1208.5846.

\end{thebibliography}
\end{document}
